\pdfinfoomitdate=1
\pdftrailerid{}
\pdfsuppressptexinfo=-1
\documentclass[11pt]{article}
\usepackage[margin=1in]{geometry}
\usepackage[T1]{fontenc}
\usepackage{lmodern}
\usepackage{amsmath,amssymb,amsthm,mathtools}
\usepackage{booktabs,array,longtable}
\usepackage[hidelinks]{hyperref}
\usepackage{microtype,enumitem,needspace}
\emergencystretch=2em
\setcounter{tocdepth}{1}
\setlist{itemsep=3pt,topsep=5pt}
\newtheorem{theorem}{Theorem}[section]
\newtheorem{lemma}[theorem]{Lemma}
\newtheorem{proposition}[theorem]{Proposition}
\newtheorem{corollary}[theorem]{Corollary}
\theoremstyle{definition}
\newtheorem{definition}[theorem]{Definition}
\theoremstyle{remark}
\newtheorem{remark}[theorem]{Remark}
\newcommand{\E}{\mathbb E}
\newcommand{\Prb}{\mathbb P}
\newcommand{\Pois}{\operatorname{Pois}}
\newcommand{\Bin}{\operatorname{Bin}}
\newcommand{\pos}[1]{\left(#1\right)_+}
\newcommand{\lam}{\overline\lambda}
\newcommand{\Exc}{\mathcal E}
\newcommand{\ind}{\mathbf 1}
\newcommand{\Wfun}{\mathcal W}
\DeclareMathOperator{\supp}{supp}
\title{Capacity Bounds for the ORAM Rebuild}
\author{}
\date{}
\hypersetup{pdftitle={Capacity Bounds for the ORAM Rebuild},pdfauthor={}}
\begin{document}
\maketitle
\begin{abstract}
Level one builds its deduplicated incoming prefix directly into a standalone
main table. At digit one it drains into the published output tier; at larger
digits it checks one global producer bucket, completes the older A-only route,
and inserts B under the same current job family. All distributing ordinary
passes use one-page A-only DFS leaves and direct B-tier offers.
The producer event is total complete-home excess, including excess spread
across several homes. Its probability is bounded on the unconditional
ideal-word space and is unioned with A-stage, A-range and final-output events.
The report gives conditional composition, complete-population excess
comparisons, the shadow recurrence, binary-top proofs and public exposure
accounting. Lean certifies all 94 listed standard $D=8$ planner
geometries at capacities $2^{11}$ through $2^{31}$, from initialized and empty
starts. \texttt{IdealApiFamily} combines arbitrary-byte-key reference API
correctness, exact probe simulation and capacity failure at most
$h\cdot2^{-64}$ over $h\ge1$ windows of $2^{32}$ accesses. The premises are
valid initialization, public geometry, shape and schedule, and logical live
capacity on every legal observation branch. No fixed key universe or numeric
key bound is required. Independent countable client seeds may have unbounded
support and seed-dependent initialization. Observation transfer gives the same
eventwise simulator error under explicit agreement off capacity failure.
\texttt{IdealApiBounded} and \texttt{IdealSeedBounded} give the same bound
on the outer measure of incorrect actual API results under explicit agreement
of actual history, issued requests and selected results. Empty-start correctness requires
empty initial history. The fixed-universe \texttt{ConfigurationFamilyWindows},
\texttt{ConfigurationFamilyOblivious} and \texttt{geometry31} statements keep
their finite-universe scope and premises.
Native certification remains open: the model's theorem does not establish
native execution or hash correspondence, caller promises or inventory-rule
semantics. Universal membership of native planner outputs in the Lean family
under the standard family settings also remains unproved. The native numerical checker therefore retains
unconditional bound one.
Exact numerical outcomes are supplied by the ledgers and
configuration-space artifacts.
\end{abstract}
\begingroup\small\tableofcontents\endgroup\clearpage

\section{Conditional composition and source contract}
\label{sec:abstract-contract}
Fix a legal logical request trace and initial data independently of ideal
routing words. For every canonical job family $f$ and distinct key $x$,
sample an independent $w$-bit word $H_{f,x}$. Copies of the same key in the
same family share one word. Potential jobs following an earlier failure
still have reference words. This section gives the fixed-trace argument.
The adaptive theorems use causal exposure bounds described in
Section~\ref{sec:lean-current}; they do not assume that adaptive requests
are independent of routing randomness.

\begin{theorem}[Conditional composition of capacity witnesses]
\label{thm:abstract-composition}
Suppose every possible first capacity failure is deterministically impossible
or implies one event $E_{e,j}$ in a finite public inventory, and
$\Prb(E_{e,j})\le b_{e,j}$ on the unconditional ideal-word space.
Suppose each component's logical-window, placement, distinctness and family
premises hold, and the exposure inventory includes initialization and jobs
crossing the horizon boundary. Then
\[
 \Prb\{\text{capacity failure}\}\le
 \min\left(1,\sum_{e,j}b_{e,j}\right).
\]
\end{theorem}
\begin{proof}
Choose the first failing test. Its pointwise implication places the
execution in the union of the inventoried reference events. Apply the union
bound. Earlier success justifies the pointwise implication; it never
conditions the probability of a reference event.
\end{proof}
No independence between jobs, A and B, or overlapping event witnesses is
required. A scheduler-dependent choice needs a union bound or a common
dominating event; the maximum of marginal probabilities alone is insufficient.
A missing event is charged one. Cryptographic and stored-key matching errors,
allocation failure and scheduler liveness require separate arguments.

\paragraph{Native source obligations.}
\label{sec:native-sources}\label{prop:native-sources}
Chronological prefix deduplication retains the newest record for each key,
including tombstones. Ordinary merges preserve the newest logical union;
adoption, forwarding, pause and lane folds preserve that population.
For a fixed trace its complete key sets are fixed independently of words.
The old and current canonical families differ at a non-top merge with $d>1$;
therefore complete old/current words are independent, although selected main
records and residuals need not be. The native audit checks public wiring and
shape compatibility. The complete history induction and applicability of
imported probability comparisons remain explicit premises of certification.

\paragraph{Unused bits.}
\label{prop:native-frontiers}
Binary routing decisions and survivor selection may inspect only the current
bits already routed. Deferred bits retain their per-key law in an A-only
frontier, before selection against B. Direct B selection at level one uses
the low $\log_2U$ current bits. This is not an independent-family B residual:
its producer and final insertion use the same job word. Neither A's reference
bound nor final-union excess is conditioned on this producer succeeding.

\section{Direct dense level one}
\label{sec:level-one-completion}\label{sec:worked-configuration}
Let $M=r_0s$ prefix buckets have $C$ slots each and $U=2M$ main homes.
After chronological deduplication the incoming population $K_B$ is distinct
and has $|K_B|\le CM=CU/2$. Direct selection scans every prefix bucket at
each public B home, retains up to $C$ winners and preserves the rest in
remaining masks. Let
\[
 X_i=\#\{x\in K_B:H_{f,x}\bmod U=i\},\qquad
 E_B=\sum_{i=0}^{U-1}(X_i-C)_+.
\]
The producer occupancy is exactly $E_B$, across all homes together.
At digit one B uses the published main/output tier directly, with capacity
$S_{\rm out}$. There is no A route, final insertion or separate narrower
producer check. At $d>1$ B uses a separate $U$-bucket main table and a
single global $C$-slot residual bucket. Its producer failure is $E_B>C$.
A failure distributed over several homes is included in this event.

\begin{proposition}[Direct B producer bound]
\label{prop:dense-producer}
Under the fixed logical B-set and ideal-family premises, the probability
of a producer overflow at capacity $S$ is at most
$\min(1,\Phi_K(U,C,C/2,S))$, where $K>S$ and $\Phi_K$ is
Equation~\eqref{eq:general-series}. Use $S=S_{\rm out}$ at digit one and
$S=C$ above digit one. Deduplication can only decrease the population.
\end{proposition}
\begin{proof}
The low-word home is exactly uniform because $U$ is a power of two.
Distinct complete B keys are independent categorical trials, so
Proposition~\ref{prop:general-tiers} applies to their full-home excess.
Selection fills each home to $\min(X_i,C)$ and the checked drain appends
each remaining winner once. Thus every drain prefix is at most $E_B$.
\end{proof}
For $C=64$, $U=16$ and $S=64$, the outward-rounded positive-series
evaluation gives about $2^{-117.72}$ per build. The conservative
$\lfloor H/(CM)\rfloor+1$ envelope at $H=2^{32}$ gives about
$2^{-94.72}$ for this component. These numbers alone are not a hierarchy
certificate; each digit has its own public exposure count in the ledger.

\paragraph{Same-family final insertion.}
Write $d=r2^e$, $U=2^k$, and
$h_f(x)=(H_{f,x}\bmod2^{k+e})+2^{k+e}
(\lfloor H_{f,x}/2^{k+e}\rfloor\bmod r)$.
Then $h_f(x)\bmod U=H_{f,x}\bmod U$.
A destination reads B main bucket $h_f(x)\bmod U$ and the global B tier;
retagging supplies the remaining direction bits. Choosing another family
for the producer would invalidate this congruence. B's scratch table must
be separate from the output A occupies before insertion.

\begin{proposition}[Composition for direct level one]
\label{prop:level-one-composition}\label{cor:direct-plan-level-one}
At digit one the per-job bound is the producer/output bound above.
At $d>1$ assume the completed-A source/frontier contract, separated offer
and drain contract, and current distinct-union law. Let $b_{A,\rm stage}$
bound the A-stage storage events, $b_{A,\rm range}$ its completed-A
coarse-range event at leaf capacity $C$, and $b_{\rm out}$ the full-union
excess event at output capacity $S_{\rm out}$. Suppose insertion stash
and lane copies have capacity at least $S_{\rm out}$. Then
\[
 b_{1,d}\le\min(1,b_{B,\rm producer}+b_{A,\rm stage}
                         +b_{A,\rm range}+b_{\rm out}).
\]
A native direct A-pool plan has $b_{A,\rm range}=0$.
\end{proposition}
\begin{proof}
Before a first failure the standalone B tier fits its producer capacity.
Complete A's route with the B region blank. Only A's old tier and all
A losses enter the one-page distributor. Proposition~\ref{prop:separated-capacity}
covers DFS storage and direct B-tier offers. Fixed insertion deduplicates
meeting main copies; combined source-rank dedup, A-before-B offers and
original-source drain dedup preserve freshness. In particular, the drain
epilogue discards an older A residual when its newer B-main copy shed into
the merge stash.
The coordinated drain of Section~\ref{sec:final-drain-bound} bounds every
insertion-stash and publication prefix by the distinct union's complete-home
excess. Since the insertion stash is at least as wide as the output tier,
one output-capacity event dominates both. Union the unconditional events;
never condition the A or union law on successful B selection.
\end{proof}
The finite-word radix density factor applies once to each appropriate
old/current witness. The producer is exact binary uniform. Copies, pauses,
lanes and repeated direct B-tier scans introduce no new random exposure.

\subsection{Reference numerical ledgers for direct level one}
Two capacity-$2^{27}$ reference ledgers are numerical evidence and inputs
to shared certificate generation; they are not an adaptive instance.
For capacity $N=2^{27}$, horizon $H=2^{32}$, $r_0=8$, $D=8$, $s=1$
and $q=2$, the ledgers give the following approximate base-two
logarithms of their exact rational conditional upper bounds:
\begin{center}
\begin{tabular}{lrr}
\toprule
Geometry & Level-one lifetime subtotal & Complete lifetime ledger\\\midrule
Uniform fanout eight & $-83.647869$ & $-69.678562$\\
Default half-fanout ladder & $-88.132529$ & $-69.678190$\\
\bottomrule
\end{tabular}
\end{center}
The level-one subtotals include producer, completed-A route/range and output
excess events with their native exposure counts. The exact values are in
\path{code/level1_certificates.json}\footnote{Paths under \path{lean/},
\path{src/} and \path{tests/} are relative to the repository root; the other
paths in this report are relative to \path{analysis/bound_certification}.} and the selected ledgers
\path{code/selected_configuration_check.json} and
\path{code/selected_configuration_half_fanout_check.json}.
Both complete ledgers meet the conditional $2^{-64}$ target with no coverage
gap. They retain status \texttt{unproved\_applicability}, unconditional upper
bound one and \texttt{complete\_configuration\_certified=false}; these
numerical results do not discharge the source or execution premises.

\subsection{Lean result and its boundary}
\label{sec:lean-current}
The numerical family covers all 94 listed standard geometries at capacities
$2^{11}$ through $2^{31}$, with bucket capacity 64, load 32, $D=8$,
automatic per-level tier/stash sizing and the $B=8$, $s=1$, $q=2$ candidate
domain. Its compact public parameters are in
\path{lean/RebuildBound/Spec/ConfigurationFamilyData.lean}; one Lean
constructor derives each member's direct, separated and top job shapes.
Arbitrary top-digit or tier pins are outside this numerical claim.

\paragraph{Arbitrary-key API endpoints.}
The canonical statement is \texttt{IdealApiFamily}: reference API correctness,
exact probe simulation and the $h\,2^{-64}$ capacity budget for legal adaptive
clients with arbitrary byte-string keys and bounded live support.
\texttt{IdealApiOblivious} transfers the budget to observations agreeing off
failure, and \texttt{IdealApiBounded} transfers it to incorrect actual API
results under explicit adapter agreement. Their independent countable-seed
forms permit seed-dependent initialization without a common key dictionary.
\texttt{ApiLegalStructural} and \texttt{IdealApiExact} expose correctness and
exact simulation under their smaller assumptions. \texttt{IdealGeometry}
provides the generic analytic/exposure forms; the numerical family specializes
the same chain. These results and the remaining native obligations are detailed
below and in \path{LEAN_CLAIMS.md}.

The source-review headline (generated by \path{lean/scripts/review_surface.py})
contains only these direct arbitrary-key interfaces. Its declaration closure
excludes the fixed key universe, numeric key encoding, finite family tape and
numerical-certificate modules. The generator exports \texttt{headlines.tsv}
separately from the complete guide and rejects fixed-universe dependencies in
the headline. This is statement navigation, not a proof.

\paragraph{Fixed-universe family construction.}
\texttt{ConfigurationFamilyStructural} checks the structural contract for
every member from both initialized and empty starts. Shared local source laws
and causal first-failure accounting feed \texttt{ConfigurationFamilyAnalytic}.
\texttt{ConfigurationFamilyNumerics} kernel-checks full-shape certificate
coverage and the exact initialized exposure sum. A symbolic identity removes
the initialization term for the empty start; no second numerical certificate
is required. The \texttt{ConfigurationFamilyWindows} theorem bounds
fixed and adaptive reference failure by $h\,2^{-64}$ over the first $h\ge1$
windows of $2^{32}$ accesses. \texttt{ConfigurationFamilyOblivious} gives the
same bound on scalar, batch and adaptive simulator distance. Statements are in
\path{lean/RebuildBound/Spec/ConfigurationFamilyCertification.lean}.

These results retain public geometry, shape and schedule, joint ideal routing
and dummy tapes, the trace and unshadowed premises, and observation agreement
off failure. For each member, one fixed universe of at most its capacity in
distinct keys with \texttt{key.toNat}$<2^{64}$ covers the whole run and every adaptive branch. The bound
is unconditional first-failure accounting, not a probability conditioned on
surviving previous windows. Exact native shape and class-clock comparisons
cover every member; they do not prove universal planner membership or native
execution refinement. Python's conditional ledger and Lean's conservative
numerical allocations remain separate artifacts.

\paragraph{Fixed-universe reference semantics.}
\texttt{HierarchyLookup} and \texttt{HierarchyUniverseLookup} establish scalar
and batch reference lookup; \texttt{HierarchyFixedSemantics},
\texttt{HierarchyUniverseSemantics} and \texttt{HierarchyAdaptiveSemantics}
establish the exact ideal-tape laws and observation transfer with error equal
to the actual failure probability. \texttt{HierarchyBatchApiCorrectness} and
\texttt{HierarchyAdaptiveApiCorrectness} establish reference API correctness
from local construction rules. These results require the structural contract,
run and schedule premises, with empty initial history for empty-start
correctness. They apply beyond the numerical family and import no numerical
matrices; other geometries still require their own numerical certificates.
\path{LEAN_CLAIMS.md} lists these exports with their premises.

\subsubsection{The geometry31 statements}
The following details describe the \texttt{geometry31} endpoints.
Their non-top fanouts are $(4,4,8,8,8,8,8)$, with $D=8$, $U_1=16$,
$T=2^{26}$, bucket capacity 64 and load 32. Published tier widths are
$(1,1,1,1,2,2,2,3)$ and stage stash widths are $(1,1,1,1,2,2,2,4)$.
Initialization drains three buckets; preparation drains four.
A workload using fewer keys, including at most $2^{27}$, is covered by these
endpoints when executed at this geometry. Smaller standard geometries are
covered by the family theorems above.

\paragraph{Deterministic reference and local laws.}
\texttt{DirectJob} assembles all-digit publication correctness, agreement
between bounded and analytical runs, and first-failure containment.
\texttt{DirectJobResume} proves fuel splitting and finite completion for fixed
geometry, inputs and words. Native interruption, helper scheduling and elapsed
time are outside this logical cursor theorem. \texttt{DirectJobReference}
extends histories to every outcome: exact published lanes on success,
otherwise empty main buckets and an unrestricted tier of newest records.
This continuation makes the reference interaction total after overflow.
Population and locality lemmas derive chronology, fixed key sets and loads.
The serialised hierarchy in \path{lean/RebuildBound/Spec/Model/Hierarchy.lean}
then supplies separated and top inputs. The geometry-specific structural,
covering and local-law obligations are discharged by
\texttt{HierarchyInvariant31}, \texttt{HierarchyCovering31} and the
\texttt{HierarchyLaws31} package. All 38 separated rows and the distinct top,
preparation and initialization bounds are kernel certified; the direct
geometries reuse existing local proofs. Shared A/B keys have one current word;
local laws are not conditioned on preceding success.

\paragraph{Fixed-trace capacity.}
The theorem \texttt{HierarchyUniform31} proves a first-horizon bound $2^{-64}$;
\texttt{HierarchyChunk31} proves that the first capacity failure lies in any
fixed chunk of $2^{32}$ accesses with probability at most $2^{-64}$; and
\texttt{HierarchyWindows31} bounds failure within the first $h\ge1$ chunks
by $h\,2^{-64}$. Statements are in
\path{lean/RebuildBound/Spec/HierarchyUniform31.lean} and
\path{lean/RebuildBound/Spec/HierarchyChunk31.lean}.

These statements fix the trace independently of routing words. Each unit has
distinct keys; each level-one ordinal has at most 512 records; initialization
keys are distinct; keys lie below $N\le2^{31}$; and the unshadowed-tombstone
promise holds under every tape. The measurable ideal family tape has one
independent uniform 64-bit word per family coordinate and key. Premises hold
at horizon $(k+1)2^{32}$ for chunk $k$ and $h2^{32}$ for the window bound.
The failure event filters potential jobs by \texttt{WithinHorizon}: a
speculative top finish consuming a carry completed after the horizon is
excluded. Windows are measured by the level-one ordinal triggering the job,
one speculative ordinal included, rather than native execution time.
The per-chunk theorem concerns the first failure in that chunk; the window
theorem unions the chunk events, without independence between jobs or
conditioning on prior survival. These \texttt{*31} statements start
from direct initialization. The family theorems above additionally cover the
reference empty-top start with zero seeds and no initialization job.

\paragraph{Encoded keys in a fixed universe.}
The generic endpoints choose one finite universe $U$ independently of the
tapes, with $|U|\le2^{31}$ and each exact key encoding \texttt{key.toNat}$<2^{64}$. Every key used anywhere
in the run belongs to $U$: initialization, queries, supplied records,
tombstones and duplicate lanes. One universe covers all chunks. These fixed-universe
endpoints do not cover unbounded churn; the canonical \texttt{IdealCapacity}
theorem below instead uses logical live capacity.
\path{lean/RebuildBound/Spec/KeyUniverse.lean} defines membership and tape
premises; \path{lean/RebuildBound/Spec/KeyRelabel.lean} states the endpoints.
\texttt{GenericHierarchyChunk31} and \texttt{GenericHierarchyWindows31}
retain the same failure budgets. One fixed key permutation transports
execution, failure events, chronological histories and probe addresses,
including after overflow, and preserves the joint uniform family/dummy law.

\paragraph{Fixed-run obliviousness and probe correctness.}
The scalar theorem \texttt{OnlineOblivious31} and batch theorem
\texttt{BatchOblivious31}, and their
\texttt{Generic} counterparts, give simulator distance at most $h\,2^{-64}$
under the joint ideal tapes and observation agreement off failure. The public
simulator sees the schedule, probing kinds and batch widths. The exact
simulator law holds on every outcome because the reference continues after
failure. Batch horizons count slots. \texttt{OnlineCorrectness31} and
\texttt{BatchCorrectness31}, and their \texttt{Generic} counterparts, establish
latest preceding supplied-value lookup, tombstone handling, earlier-lane
forwarding and address agreement for arbitrary routing and dummy words,
without a capacity-success assumption. Their bounded-key statements are in
\path{lean/RebuildBound/Spec/Oblivious31.lean}.

\paragraph{Adaptive security.}
\texttt{AdaptiveExact31}, \texttt{AdaptiveWindows31} and
\texttt{AdaptiveOblivious31}, in \path{lean/RebuildBound/Spec/Adaptive.lean},
use the same geometry and budget. A deterministic client chooses each
operation's records from the rows of all earlier operations. Batch lanes are
chosen together. The public shape (probing kinds and batch widths) and plan
schedule are fixed by the program before the tapes. On every candidate
transcript the replay must satisfy \texttt{AdaptiveSetting}, including one
fixed universe $U$ of at most $2^{31}$ distinct keys with \texttt{key.toNat}$<2^{64}$ across all
branches, the structural setting and unshadowed promise.
The transcript is measurable and has exactly the existing public simulator
law. Its adaptive failure event over $h\ge1$ windows has probability at most
$h\,2^{-64}$. Any measurable observation agreeing with a measurable image of
the transcript off that event is within $h\,2^{-64}$ of the simulator on
every measurable event. \texttt{AdaptiveTwoClients31} bounds two such clients
with the same public shape, schedule and observation map by $2h\,2^{-64}$.

The proof uses causal exposure: a family's population is complete before its
words are observed. Separated A witnesses are bounded at the early population
cutoff and B/union witnesses at the late cutoff. It therefore permits B to
depend on observations of A's old family. \texttt{BatchFound31} and
\texttt{BatchFoundCausal} establish word-free, causal selected records;
\texttt{FlagSoundness31} and \texttt{AdaptiveFlagAdmissible31} derive the
unshadowed promise for clients following the model's flag rule.
The native application still needs causal public observation maps, fixed
before the tapes and shared by the compared secrets, and agreement of the
engine's selected records and observations with the model off failure.
\path{LEAN_CLAIMS.md} (``Adaptive clients'') states that boundary.

\paragraph{Reference API semantics and native boundary.}
The API correctness theorems are \texttt{BatchApiCorrectness31} for fixed runs
and \texttt{AdaptiveApiCorrectness31} for adaptive runs. They establish Read,
Write, Erase, Exchange and WriteWithoutProbe logical-map transitions,
including prior-value responses, duplicate batch lanes and present zero
values. Local native-adapter and access-lowering theorems make their premises
explicit. Native execution correspondence, worker scheduling, compiled-code
timing and the AES-to-ideal reduction remain outside the proof.
Rebuild memory traces are assumed to depend only on public geometry.

Public internal chunks stop at requester-unit boundaries.
\texttt{ApiGroupedChoices} proves that repeated internal adapter calls retain
the requests chosen from the original pre-batch observation prefix; arbitrary
later observations do not create extra client decision points. The resulting
shape and slot schedule depend only on public descriptors and geometry.
\texttt{ApiGroupedReference} proves exact chronological record composition
across these cuts, including encoded tombstone flags.
\texttt{ApiExecutionImplements} constructs physical postrecords for fixed
request groups using provisional lane probes; causality proves their local
selection equations without a universe premise, and their flags follow.
\texttt{ApiExecutionCorrectness} derives actual-tape API correctness from the
existing explicit fixed-universe premises. \texttt{ApiAdaptiveReplay} and
\texttt{ApiAdaptiveInteraction} connect the request client to replay and
actual execution, using one observation prefix per public batch.
\texttt{ApiAdaptiveCorrectness} derives the adaptive setting and API results
under the same explicit universe premises. Physical selections need not equal
chronological reference selections when publication has discarded a tombstone.
The registered byte-key adapters and canonical product-law package prove
length-sensitive identity, finite uniform restrictions, independent blocks
and sectionwise bounds. \texttt{RequesterCuts} proves complete frozen-unit
and grouped trace-prefix chronology, including cuts in the finite trace's
empty tail. Such cuts saturate at the last frozen unit, excluding a partial
current unit. \texttt{HierarchyStructural} proves top
history and finish structure without a public numeric key bound, retaining
explicit initial length, input-rate and local flag premises. \texttt{ApiCapacity}
counts exact logical support at every lane prefix, including unbounded misses,
churn and capacity-preserving idle extensions. \texttt{HierarchyLivePopulations}
derives the four top-job populations for every modeled top ordinal from that
support, without a supplied-cut premise,
conditional on named API-prefix correspondence and flag premises. The
\texttt{ApiStructural} package derives both premises from physical request
construction, for fixed and adaptive runs, without a key universe or global
key-count bound. Deterministic correctness needs only actual initialization
size; live populations additionally need logical prefix capacity and unit-safe
public grouping. Empty future trace inputs do not append physical operations.
The \texttt{ApiContinuation} package constructs idle physical padding to a
complete-unit endpoint, preserving prior records, selections and bucket
observations. Its completed trace history represents the original final map;
extended horizon and schedule validity remain explicit. \texttt{ApiLegal}
provides public, slot-dependent bucket ranges and typed client prefixes, with
capacity required only on legal transcripts; actual observations are unchanged
and public batch decision points are preserved. \texttt{IdealHierarchy} feeds
the hierarchy from the canonical routing/dummy tape, proving carry and horizon
coherence, finite joint laws and outside-dependent finite-key section bounds.
\texttt{IdealMeasurable} and \texttt{IdealFailure} prove measurability of
actual adaptive transcripts and capacity-failure events on that product tape,
including reference continuations after overflow. Finite local word dependence
follows the actual computation and requires no shared key dictionary or success
premise. \texttt{IdealSeed} proves independent-seed mixture and unconditional
coupling adapters; \texttt{ApiSeed} applies them to legal clients with
seed-dependent initialization and countable seeds of unbounded support.
Seedwise exact laws and error estimates remain named premises of the adapters.
The actual adaptive input prefix has canonical-tape exposure causality
and independence from the entire fresh routing family
(\texttt{IdealExposure}, \texttt{IdealFreshness}). Initialization has its actual
analytic failure bound from the supplied input count
(\texttt{IdealInitialization}), with no global key universe.
This does not make routed intermediate layouts independent of their words;
their local laws must account for that dependence.
\texttt{IdealSepUnion} bounds the actual separated current-home union term
for all certified geometries and both starts. Exposure fixes the complete
input key sets while current-family routing may change their physical layout;
requester widths supply the populations without a global key universe.
\texttt{IdealSepA} derives the separated early-A route and range bounds by
sampling the old and current families jointly. \texttt{IdealSepFirstFailure}
combines those witnesses with the union term for all 94 geometries and both
starts, without a global key universe or an earlier-success assumption on the
reference execution. A private finite restriction supports fixed-universe lower-table
locality, without restricting the actual adaptive keys.
\texttt{IdealTopPreparation} derives the top-preparation contribution
from legal adaptive API clients and logical live capacity. The two-family
section keeps the previous top's layout dependent on its old words and proves
the exact binary old-home density. \texttt{IdealTopPreparationFirstFailure}
specializes the existing analytic budget to all 94 geometries and both starts.
Initialization, public width/schedule and branchwise live capacity are explicit;
no per-job population or key-universe premise is supplied. Tier fit is part of
the event, without conditioning the reference execution on earlier success.
\texttt{IdealTopFinish} derives the corresponding finish contribution from
the same legal API contract. Its single-family section preserves complete
record sets while preparation and the completing carry share the current
family. Shared A/B keys reuse a word; other source keys supply independent
finite history. Logical capacity bounds both A and resolved live winners.
\texttt{IdealTopFinishFirstFailure} applies the existing analytic budget at
cutoffs 576 to all 94 geometries and both starts.
Direct jobs have canonical-tape bounds for B, union, and the two-family
old-home A witnesses. \texttt{IdealDirectFirstFailure} composes them for all
94 certified geometries and both starts, including reference continuations
after overflow, without a key universe or supplied probability premise.
Canonical records, requests, API maps and routing inputs use exact finite
byte strings, preserving length and every byte. Numeric fixtures and
finite tape views use explicit exact encodings. \texttt{IdealProbe} proves
exact reference probe simulation without a live-capacity premise, using a
separate finite restriction for each replay fibre. \texttt{IdealCapacity}
composes all first-failure classes with the existing exposure counts and
certificates, giving $h\,2^{-64}$ for all 94 geometries and both starts over
$h\geq1$ windows of $2^{32}$ accesses. \texttt{IdealApiFamily} combines
request correctness, exact law and that budget from valid initialization,
public execution premises and logical capacity on every legal observation
branch. Independent countable client seeds may have unbounded support and
seed-dependent initialization. \texttt{IdealApiOblivious} and
\texttt{IdealSeedOblivious} give the same eventwise simulator error under
explicit agreement with the reference off failure.
\texttt{ApiBoundedTransfer} supplies the generic arbitrary-key bounded API
reduction under structural geometry, actual initialization size and public
schedules, without live capacity or a tape-distribution assumption. Actual
history, issued requests and selected probe results must agree with the
reference off its capacity-failure event. The conclusion covers every logical
prefix, scalar and batch replies, and probing addresses. Requiring the actual
request count to equal the public slot count and actual slot count excludes
truncated executions. \texttt{IdealApiBounded} and
\texttt{IdealSeedBounded} compose this containment with the certified capacity
budget for all 94 geometries and both starts. Incorrect-result outer measure
is at most $h\,2^{-64}$ over $h\geq1$ windows of $2^{32}$ accesses; adapter
measurability is not required. The seed theorem permits independent countable
seeds with unbounded support and seed-dependent initialization. No event is
conditioned on success. Concrete native agreement, cryptographic error
accounting remain separate. The fixed-universe probability endpoints
retain their finite-universe and encoding premises.

\texttt{IdealGeometry} exports the generic canonical analytic chain.
Explicit geometry, width, tier-link and cutoff checks derive all actual local
probabilities; initialization and logical live capacity derive populations.
A public exposure-sum inequality supplies the error, with optional window
arithmetic yielding $h$ times a base-window budget. Reference correctness,
exact simulation, independent countable-seed capacity and off-failure
observation transfer are also exported. Load 32, bucket width 64, fanouts
2--8, depth at most 14, lower cutoff 256 and compatible direct/DFS shapes
remain explicit restrictions; top cutoffs are parameters.
\texttt{IdealGeometryFamily} proves all 94 geometries and both starts
instantiate that interface with exactly the existing class budgets.
\texttt{IdealCapacityAnalytic} specializes the generic window proof.
This changes neither the numerical family nor the native error obligations.

Finite differential replays (\path{lean/oracle/}) compare engine and model
tables lane by lane on small configurations and on two capacity-$2^{27}$
reference configurations, where only lower levels are evaluated. This evidence
supports the shared implementation, but does not establish an adaptive
instance or a universal native refinement.
\path{LEAN_CLAIMS.md} states the claims, premises and evidence.

\section{A simpler distributor: containment of range events}
\label{sec:dfs-theorem}
If each coarse destination range fits, a parent buffer also fits when it has
room for all its coarse descendants. Finer buffers contain subsets of one
coarse range. This containment lets one probability bound cover the entire
distributor. Depth-first traversal keeps only a path and pending siblings
in memory, while still visiting every public node.

The reference construction first expands the padded distribution tree and
then performs the existing merge and residual offers. DFS is an evaluation
schedule for that construction: preserve the same local operations and
dependencies, but materialize only buffers still needed. It saves memory
and can avoid memory passes; traversal itself adds no randomness assumption
or new deduplication rule. The node-capacity argument below applies to either
evaluation schedule. Equivalence here concerns this expanded padded tree.

\begin{proposition}[Public-capacity DFS containment]
Let a swept residual pool have at most $P_{\rm pool}$ physical records, and partition
its destination domain into $Q=2^{D_c}$ public coarse ranges. Unreachable padded
ranges contain no records. Suppose a first-failure reference event
\[
 E_c=\{\text{some coarse range contains more than }S\text{ candidates}\}
\]
has a proved probability bound $\beta_c$. Use a full binary depth-first traversal
of the coarse partition, with any fixed number $R_f\ge0$ of additional
refinement levels. Give every depth-$j$ node public capacity
\begin{equation}\label{eq:dfs-capacity}
 C_j=\min\bigl(P_{\rm pool},S2^{\max(D_c-j,0)}\bigr),\qquad 0\le j\le D_c+R_f.
\end{equation}
Assume lossless oblivious partitioning into padded children, with work arrays
large enough for the entire parent input, a checked child capacity, and a
public visit schedule. Then any first distributor overflow implies $E_c$.
Its probability is at most $\beta_c$, with no extra factor for nodes or depth.
\end{proposition}
\begin{proof}
On $E_c^c$, a depth-$j$ node above the coarse frontier contains at most
$2^{D_c-j}$ coarse populations of at most $S$ records each, and never more
than the entire pool's $P_{\rm pool}$ records. At the frontier its demand is at most
$S$. Every finer descendant is a subset of one coarse population, so also
has demand at most $S$. Thus every node fits Equation~\eqref{eq:dfs-capacity}.
Conversely, an overfull coarse population cannot pass its frontier capacity
test in a lossless execution. If gathering the root requires compaction
from $P_{\rm pool}$ input slots to $C_0<P_{\rm pool}$, its failure has the same containment.
For an ideal reference history with an earlier failure, retain the original
coarse theorem's first-failure witness; no conditioning on earlier success
or fresh split-bit law is introduced.
\end{proof}

The partitioner may use stable compaction separately for the two bit values,
followed by fixed-size padded copies. Its temporary arrays must hold the
whole input; using another capacity-limited butterfly inside the split would
reintroduce a missing event. Visit both children, including empty subtrees,
without secret-count-dependent addresses or allocation. Source identities,
tags, masks and freshness metadata must be preserved.

The native distributor takes $R_f=0$: its leaves are the coarse ranges.
Depths whose capacity is the whole pool cannot fail, so it compacts the
gathered pool directly into each node of the first binding depth, then
halves capacities by binary splits, compacting each right child in place in
its dead parent's buffer. Appendix~\ref{sec:dfs-application} gives its
schedule and storage. These DFS symbols are separate from the merge digit
$d$, odd radix $r$, binary width $b$ and early-source cap $P$ in the
population calculation.

\paragraph{Selected two-source specialization.}
Take the pool to contain only A's old tier and A-only route residuals, set
$S=C=64$ and $R_f=0$, and scan B's existing $s\ge1$ global tier pages
directly for every destination. Combined swept-source deduplication still
precedes distribution; splitting storage does not split that operation.
The following proposition is the proof contract for this design.

\begin{proposition}[One A leaf and global B pages]
\label{prop:separated-capacity}
Consider a pass with exactly two swept sources: A's residuals and a
completed B tier of public capacity $Cs$. Let $A_R$ count the A residuals
in coarse range $R$ before any deletion against B, with reference tails
$\Prb\{A_R>C\}\le a_R$. Assume public lossless DFS capacities as in
Equation~\eqref{eq:dfs-capacity} with $S=C$, sufficient work arrays, and
no replication of B's payload per destination. Preserve combined source-rank
deduplication, original-source lifetime, freshness-aware A-before-B offers,
full-home matching for B, private consumption masks and their lane folds,
and the coordinated epilogue. Then:
\begin{enumerate}
\item Any first A-distributor overflow implies
      $E_A=\{\max_R A_R>C\}$, with probability at most
      $\min(1,\sum_R a_R)$.
\item Every destination scans one A page and the same $s$ B pages,
      at most $C(1+s)$ physical candidate slots. B adds no distributor
      event; its tier-capacity failure is charged at its producer only.
\item Under the ordinary final-route premises, the combined leftovers
      and final sheds retain the complete-home excess witness
      $E_{\rm home}=\sum_v(X_v-C)_+$ for the two drain phases.
\end{enumerate}
No independence between A and B, or conditioning on successful routing,
is required.
\end{proposition}
\begin{proof}
Combined deduplication only deletes A records, so the routed A count in
$R$ is at most $A_R$. On $E_A^c$, DFS containment bounds the gather and
every node by its public capacity. B already occupies at most $Cs$ slots
before the first failure and is read directly, so the offered slot count
is deterministic; no mixed range is materialized or capacity-tested.
The union bound gives the stated A charge on the unconditioned reference
space. For draining, combined dedup leaves at most one swept copy per key.
Freshness-aware offers never reduce main occupancy and leave a candidate
only at a full home. Fold consumption masks and run drain-only dedup
against the original fresh sources, including consumed records. Drain the
A leftovers and surviving B tier once into one output, then perform the
existing leftover dedup and append surviving final sheds. The argument in
Appendix~\ref{sec:final-drain-bound} therefore bounds every packing prefix
by $E_{\rm home}$. Earlier B-producer failure is covered by its original
witness, not charged again at each scan, range, lane or consumer.
\end{proof}

Thus the new distributor term in the lifetime ledger is simply
$\sum_j\sum_R a_{j,R}$ over eligible public pass ordinals $j$.
Preparation stashes, B's producer tier, final sheds and output packing keep
their existing charges. Section~\ref{sec:completed-a-shadow} supplies the completed-A marginal.

\paragraph{The same pass at the top.}
The top finish has the same two swept classes. A is the preparation's
retained carry R: A's route and offers complete in the top's preparation,
so R holds exactly the records of A's full homes that did not fit. B is
the incoming carry's tier, offered directly. Only the drain witness
differs: the top drops every tombstone, and its merge stash and output
tier are bounded by the excess of the resolved live set rather than of the
merged union. Corollary~\ref{cor:top-finish-separated}
(Appendix~\ref{sec:top-excess}) states the top finish as this instance,
with A marginal $a_R=\Phi_q(2^I,C,32,C)$ for coarse ranges of $2^I$ homes
and the drain witness of Proposition~\ref{prop:top-excess}. No pass of the
level arena routes a range table any more, and the top digit $D$ enters the
top finish's ledger only through its exposure count.

\section{Notation and scope}
The probability bounds are functions of public configuration $\Theta$ and
access horizon $H$. Legal native geometry, applicability of a probability
law, and meeting the requested target are separate checks.
\begin{center}\small
\begin{tabular}{p{0.22\linewidth}p{0.69\linewidth}}
\toprule
Public parameter & Meaning\\
\midrule
$L,\ b_\ell,\ U_\ell$ & Hierarchy depth, fanout and carry width at each level.\\
$C$ & Main-bucket record capacity.\\
$S_e$ & Capacity of the particular stash, range bucket or tier being tested.\\
$N_i,\ m_i$ & Complete logical population bound and home count of source $i$.\\
$\lambda$ & An upper bound on the relevant mean loads $N_i/m_i$.\\
$w,\ \mathcal R$ & Routing-word width and public destination/scatter maps.\\
$H,\ \varepsilon$ & Access horizon and requested lifetime probability target.\\
\bottomrule
\end{tabular}
\end{center}

Family identities, ownership, scheduling and deduplication order are also
part of the model. Symbolic $C$ in a theorem does not imply a configurable
native SIMD width. In the worked ordinary example, use:
\begin{center}\small
\begin{tabular}{p{0.18\linewidth}p{0.73\linewidth}}
\toprule
Symbol & Meaning\\
\midrule
$\ell,d$ & Ordinary level and public merge digit, $1\le\ell<L$, $1\le d\le b_\ell$.\\
$U=U_\ell$ & Incoming carry width, $16\prod_{i<\ell}b_i$ buckets: $16\cdot8^{\ell-1}$
on the uniform ladder; 16, 64, 256, 1024, 8192, 65536 on the default one.\\
$A,B$ & Older partial table and incoming table; they may share keys.\\
$M=(d-1)U$ & Older main width for $d>1$; output width is $dU$.\\
$h_{\rm old},h_{\rm new}$ & A's previous routing family and this merge's family; B already uses the latter.\\
$X_i,a_i$ & Complete logical old-home count and retained main count at bucket $i$.\\
$S$ & One physical stage stash's capacity, in records.\\
$D$ & Upper likelihood ratio for the old map relative to uniform old homes.\\
$I$ & Low destination bits left for a main/range offer to match.\\
\bottomrule
\end{tabular}
\end{center}

\section{Abstract history and storage lemmas}
\label{sec:abstract-history}
These lemmas discharge algebraic and scheduling steps once, independently
of a particular ladder. Their hypotheses specify what remains to be proved
about native transitions. They use the fixed logical trace and ideal-word
experiment of Section~\ref{sec:abstract-contract}.

\begin{lemma}[Windows and canonical families for variable fanouts]
\label{lem:window-family}
Let $a$ be the number of access positions in one prefix, let
$b_1,\ldots,b_{L-1}\ge2$ be integer fanouts, and set
$P_\ell=\prod_{i<\ell}b_i$. The $v$th incoming carry at level $\ell$
represents the interval
\[
 W_{\ell,v}=[(v-1)aP_\ell,vaP_\ell).
\]
Assume each prefix logical set is fixed by its request subsequence, has
at most one record per position, and is distinct by key. Merges retain
exactly the newest record of each key in their input union, and forwarding,
adoption and pause preserve the represented logical set and window.
Then a non-top partial at digit $d=1+((v-1)\bmod b_\ell)$ represents
$[(v-d)aP_\ell,vaP_\ell)$ and has at most $daP_\ell$ distinct keys,
counting tombstones. At a merge with $d>1$, A and B have bounds
$(d-1)aP_\ell$ and $aP_\ell$ respectively.

Define $F(\ell,v)$ by repeatedly replacing $(\ell,v)$ with
$(\ell+1,v/b_\ell)$ when $\ell<L$ and $b_\ell$ divides $v$.
The endpoint $vaP_\ell$ is invariant under this replacement, producer
and consumer have the same canonical family, and the older and current
families at every non-top merge with $d>1$ are different. Consecutive top
ordinals also have different families.
\end{lemma}
\begin{proof}
The prefix assertion is the base case. A partial is the union of $d$
consecutive carry intervals; deduplication cannot increase their total
number of access positions. At $d=b_\ell$ their union is exactly
$W_{\ell+1,v/b_\ell}$, since $P_{\ell+1}=b_\ell P_\ell$.
This proves the window and population assertions by induction. Preserving
operations do not change them. The same product identity proves invariance of the endpoint
and producer/consumer agreement. For $d>1$, $v-1$ is not divisible by
$b_\ell$, so the older family's terminal coordinate is $(\ell,v-1)$.
The new family's terminal coordinate is $(\ell,v)$ or has a larger level;
both choices differ from the old one. At the top the coordinate is $(L,v)$.
\end{proof}
For a non-top output of $dU_\ell$ homes the implied load bound is
$aP_\ell/U_\ell$; it need not be fixed to one load at every level.
Thus changing the lower fanouts or enlarging level-one tables does not
change the proof. Initialization may seed counters at a public cycle
boundary with no lower partials; the subsequent windows are translated by
that boundary. Its initial top population needs its own bound and does not
inherit a non-top window bound.

On a fixed trace, the complete key sets specified by these windows are
fixed independently of routing. In the ideal model, distinct family
coordinates therefore give independent old/current word arrays on those
sets; a shared key has one current-family word. This is an unconditional
statement about complete logical sources. Conditioning on selected tiers,
main survivors or earlier success can destroy independence. The native obligations in Section~\ref{sec:native-sources} still require
proof that the sources implement these sets, that family encoding is
injective without wraparound (including initial families), and that the
listed transitions preserve them.
The functions \texttt{ForwardToTerminal} and \texttt{FamilyWindowEnd}
in \path{src/oram/levelarena/cpu_byteslice_level_arena_family.hpp} implement
the algebra above; finite coordinate checks alone would not establish the
source induction.

\section{Lifetime accounting}\label{sec:ledger}
Let $\mathcal W(\Theta)$ be a set of witness classes covering every possible
first capacity failure. For class $e$, let $n_e(\Theta,H)$ bound its number
of exposures, including initialization and jobs crossing the horizon
boundary. For every exposure, prove that a first failure implies a
reference-word event with probability at most $b_e(\Theta)$. The ideal
words in this argument are not conditioned on earlier capacity success.
Then the union bound gives
\begin{equation}\label{eq:general-lifetime}
 \Prb\{\text{any ideal-tag capacity failure}\}
 \le B(\Theta,H)
 :=\min\!\left(1,\sum_{e\in\mathcal W(\Theta)}
               n_e(\Theta,H)b_e(\Theta)\right).
\end{equation}
Witness events need not be independent. If several theorems bound the
same witness, their minimum may be used; they are not additional failures
to charge. Alternative public shapes within one exposure can instead be
handled by a maximum when a common exposure bound is used. Cryptographic
distinguishing and key-equality errors are separate from this ideal-tag sum.

An omitted event invalidates a complete certificate. It must be reported
as unproved, or conservatively assigned an upper bound of one. A small sum
over the currently covered events alone is not $B(\Theta,H)$.


\paragraph{Separate cryptographic and matching budget.}
Let $\delta_{\rm crypto}$ bound the advantage for the tested outcome event
between the native seeded construction and a specified ideal primitive experiment. Suppose a
coupling with the reference capacity experiment agrees through first failure
except for normalization collisions, false stored-key matches or routing-input
aliases, of unconditional ideal probabilities at most $\eta_{\rm norm}$,
$\eta_{\rm match}$ and $\eta_{\rm alias}$, or resource refusal of probability
at most $\eta_{\rm stop}$. Then
\[
 \Prb_{\rm native}\{\mathrm{failure}\}
 \le \min(1,\delta_{\rm crypto}+\eta_{\rm norm}
              +\eta_{\rm match}+\eta_{\rm alias}+\eta_{\rm stop}+B(\Theta,H)).
\]
\texttt{NativeFailureBound} checks this composition.
The \texttt{NativeObservationBound} transfer is two-sided for a tested public
event under exact reference simulation and off-failure
agreement. The eventwise hybrid remains a premise; a public-view hybrid
alone does not bound incorrect confidential API results. These statements
prove neither AES security nor native correspondence.

In the exact configuration, wide-key rebuild helpers (stored widths 9--15
always; width 16 without VP2INTERSECT or with
\texttt{CPU\_BYTESLICE\_ORAM\_WIDE\_KEY\_VP2INTERSECT} off) and the
VP2INTERSECT dense direct-slot matcher at width 16 with that option off compare both
qwords against the same partner before reducing to membership. Other dense
lookups use the generic complete-key matcher. The OR of the
two XOR differences is zero exactly for complete stored-key equality;
unsigned minimum, synchronized rotations and complete-zero repair preserve
that predicate. A checked Lean experiment proves this algebra and slot
coverage; scalar differential regressions and generated-code inspection
exercise its native implementation. Exact equality contributes zero $\eta_{\rm match}$ subject to
native correspondence, including packed-byte decoding and dense uniqueness.
It does not remove normalization collisions or establish complete execution
refinement. The default VP2INTERSECT build instead matches sixteen-byte AES
stored keys approximately: rebuilds AND four 32-bit segment memberships and
dense lookup ANDs two 64-bit halves, so different partners can supply
different segments. The registered \texttt{NativeMatch} statements prove that
this rule has no false negatives and agrees with key membership off its
false-match event. The registered \texttt{NativeMatchLanes} package proves
finite packed-word, unpack, mask-placement, zero-repair, retained-source and
class-direction semantics, with conditional bridges to exact duplicate masks.
\texttt{NativeMatchDense} proves the qword transpose, full-key verification
and coverage of a lost present hit by fourteen other row/column subsets,
under live-key uniqueness. Memory/byte layout and compiled execution remain
separate. For a comparison inventory independent of the stored keys,
they bound the false-match probability by $\min(1, \sum m^k / 2^{128})$ for
$m$ partners and $k$ segments, under a random function or an ideal
permutation. A per-site model of the public rebuild schedule counts these
comparisons. Charging real records only, it bounds the
term by $2^{-65.3}$ per $2^{32}$-access window over all certified geometries,
about $2^{-97.3}$ per access; dense lookup adds at most $2^{-80.5}$. The registered production hierarchy inventory derives the tighter
$S_{\min}$ count from actual comparison prestates and finite-prefix carry
ordinals. Exact arithmetic combines it with the existing Lean capacity
allocation for all 94 geometries and both starts, proving capacity union
matching at most $h2^{-64}$ over $h$ windows. The finite native adaptive API,
sparse routing replay and its logged local diagnostics are modeled.
The checked \texttt{NativeMatchBestCodedFailure} and
\texttt{NativeMatchBestCodedObservation} endpoints transfer correctness and
public observations to an independent ideal RF64 oracle on actual complete
native messages. For fixed certified geometry and initialization mode $g$,
$I$ initialization records and $n\le h2^{32}$ actual charged lanes, their
normalization allowance at fixed widths seventeen through sixty-four is
\[
 \eta_{\rm norm}(g,I,n)=\min\!\left\{1,
   \frac{\min\!\left\{\binom{I+n}{2},\,A(g)+nB(g)\right\}}{2^{128}}\right\}.
\]
At width sixteen, $\eta_{\rm norm}=0$ under the existing ideal-permutation
model. The birthday and affine endpoints concern the same native failure
and public-observation events, so this comparison needs no additional
error term.

The coefficients count exact stored-key comparisons and deduplicated pairs
of keys read within each complete-message routing family. For an explicit
conservative construction, let $L$ be the number of levels, $S$ the allocated
table-slot count, $r_0$ the requester units per prefix, $N$ the capacity,
$P_\ell$ the level-$\ell$ carry period in charged lanes, and $f_\ell$ its
fanout. In the certified load-32 family, put
\[
 \begin{aligned}
 W&=\max_{1\le\ell<L} f_\ell P_\ell,
 &M&=\max\{N,W\},\\
 K&=L(3M+64r_0+N)+SW+3P_1+64,
 &C&=\binom K2,\\
 A(g)&=LC,
 &B(g)&=192+192r_0+SM+\sum_{\ell=1}^{L}
     \left\lceil\frac{C}{P_\ell}\right\rceil.
 \end{aligned}
\]
Here $M$ bounds logical table populations even on reference overflow
continuations, and $K$ bounds the combined construction and probe keys of
one family on ring-ready branches. The canonical pair inventory is empty
on ring-stop branches, whose probability is charged separately. These
natural-number coefficients depend only on the fixed geometry/start, not on $h$, $n$, key identities or random outcomes. Blind
lanes count toward $n$. The normalization count uses actual $n$, not the
certificate horizon; the separate finite replay depth is not a pair-count
coefficient. The numerical report evaluates these bounds for all 94
geometries and both starts. Their conservatism does not guarantee an
improvement over the birthday bound before saturation for every geometry.

The common error is
$\varepsilon=h2^{-64}+\eta_{\rm norm}+\operatorname{ringCost}$.
Failure is at most $\min\{1,\varepsilon\}$; each observation inequality adds
that same error to the other event probability and caps at one. Exact
comparisons and routing use one shared normalization charge, additional to
the existing capacity/matching allocation. The public view is the optional
address-row transcript, absent on error or stop; confidential replies and
additional compiled leakage are excluded. Native outcomes after a bad event
remain unrestricted. The ring-stop probability and both routing marginal
laws are evaluated on unconditional experiments. This local ideal bound
does not remove the quadratic lifetime-query terms in the cryptographic
reductions below, or cover normalization calls outside the modeled run.
Stored-key secrecy, cryptographic
replacement (including the width-sixteen permutation), normalization lifetime
allowances, public encoding/counter fit, schedule realization, other
resource/completion stops, simulation time and compiled memory/concurrency/
instruction execution remain separate obligations. No total native error
bound is claimed. The hierarchy derives routing and normalization keys by
encrypting two fixed distinct labels under its master AES key. With an
independent uniform secret master used only for those labels, separation costs
the root AES-PRP advantage at two queries plus $2^{-128}$ for coupling a uniform
distinct pair to independent keys. The reduction simulates the entire
downstream execution; this cost is included once in $\delta_{\rm crypto}$.
Routing uses a fixed two-block
CBC message: the labeled packed identity followed by the complete stored key.
This excludes wide-key message aliases, subject to identity bounds
and native correspondence. The registered \texttt{NativeRouting} package
identifies the native mask/shift expression with bounded arithmetic packing
and proves complete-message injectivity, including real/dummy separation.
Discarded-bit counterexamples show that the field bounds are necessary;
the package's finite lane identities also check shared-key interleaving,
variable-message assembly and main/tail output gathers for arbitrary block
operations. Byte packing, AES rounds, native loop/memory behavior, compiled
execution and allocated-family correspondence remain explicit. The package's
family adapter derives canonical-family field bounds from the constructor's
conservative lifetime allowance, including seeded and speculative ordinals for
both starts. It reuses the hierarchy's forwarding function and preserves
producer/consumer sharing. Native geometry mapping, actual counter progression
and publication correspondence remain separate obligations.
The fixed-length CBC theorem gives routing distance
at most the AES-PRP advantage at $2q$ block queries plus $8q^2/2^{128}$,
clipped to one, for $q$ routing messages; it already includes the permutation
switch. At $q=2^{32}$ the statistical term is $2^{-61}$, so this construction
does not establish a total native $2^{-64}$ target. At each fixed native key
width $u>16$, normalization uses $m=1+\lceil u/16\rceil$ CBC blocks. The same
theorem supplies the AES-PRP advantage at $mQ$ queries plus
$2m^2Q^2/2^{128}$; in its random-function hybrid, distinct-key collisions cost
at most $Q(Q-1)/2^{129}$ separately. The checked
\texttt{NativeNormalize} package derives fresh-answer laws from replay
fibres, constructs those fibres for causal fresh-query recursion, and integrates
arbitrary independent measurable environments, including infinite oracle tapes.
The same package proves the birthday bound and constructs answer slots for
bounded-width conversion replay on the same unconditional space. Repeated keys
reuse cached answers; repeated and unused slots query distinct proof-only keys
longer than the real-key width. Returned conversions satisfy the persistent-tape
step equation, and padding answers are hidden from the chooser. A direct trace
endpoint derives replay correspondence from a chronological conversion
list, its budgeted length, and causal local equations through the unused suffix.
Unused reservations add no trace entries. Establishing these premises for
native execution and its independent environment remains an obligation;
the query limit alone does not establish them. Padding adds
no native primitive calls. Keys below sixteen bytes are stored verbatim, and
a sixteen-byte key is one AES permutation; neither needs a random
normalization replacement for collisions, though segment matching needs ideal
sixteen-byte stored keys. The checked \texttt{NativeKey} byte-encoding argument
establishes injectivity at fixed width, using injective block encryption only at
sixteen bytes. It also proves long-key input-message injectivity
and the exact CBC block count; this does not assert digest injectivity or a
mixed-width cryptographic bound. Native byte packing and AES correspondence
remain explicit. Entropy,
native correspondence and complete lifetime budgets remain open.
\path{ERROR_ACCOUNTING.md} specifies the
experiments and outstanding query-budget obligations.
The conditional capacity target alone does not allocate these extra terms.

The native engine enforces the declared online access horizon and checks
scalar/batch dummy-counter reservations before probing. Composed callbacks and
tree batches check their full access count before mutation. Native boundary
regressions check refusal, state preservation and completion of accepted work.
Construction also checks seeded positions, speculative family ordinals and
complete window endpoints over that horizon. The same bound covers padded
dense windows and shifted identifiers; hashed-table allocation stops below
the dense identifier namespace. A separate, checked normalization reservation
limit $B$ defaults to capacity plus the access horizon. It bounds both distinct
inputs and actual conversions, including standalone conversion calls and failed
attempts. Initialization and batch chunks reserve all inputs before conversion;
rollback does not refund them. At width 16 this bounds normalization AES
block calls by $B$, and above sixteen by $mB$. A separate routing allowance $R$
reserves every leaf evaluation before encryption, including padding, repeated
worker work, speculative rebuilds and failed initialization. It bounds distinct
routing messages and actual evaluations by $R$, and routing AES block calls by
$2R$. Its default is only an arithmetic ceiling. Refusal permanently stops the
instance and may interrupt accepted work; it provides no rollback or completion
guarantee. The full-run API contract requires the advertised request count,
so such refusal belongs in $\eta_{\rm stop}$. Query limits alone do not bound
its probability: use one unless a smaller bound or a matching reference stop
contract is justified. Neither allowance supplies a simulation-time bound or
establishes native execution/observation correspondence. A stopped reference
requires a separately proved simulator law and does not imply full-run API
correctness. The native error
budget remains conditional.

\section{Reusable computable probability bounds}
The following results are parameterized statements. Their applicability
depends on the source law, not on the number of rows in a geometry audit.

\subsection{An uncapped-arrival bound}
\begin{proposition}[Subset witness]\label{prop:general-subset}
Suppose a storage demand is at most $\sum_{v=1}^n(Z_v-C)_+$ and, for every
node subset $J$ and $z\ge1$,
\[
 \E z^{\sum_{v\in J}Z_v}
 \le D\exp\{\lambda |J|(z-1)\},
 \qquad 0<\lambda\le C.
\]
For any nonnegative integer storage capacity $S$, its overflow probability
is at most
\begin{equation}\label{eq:general-subset}
 \mathcal B_{\rm sub}(n,C,\lambda,S,D)
 =\min\!\left(1,D\sum_{j=1}^n\binom nj
       \frac{\exp\{\lambda j(z_j-1)\}}{z_j^{Cj+S+1}}\right),
 \qquad z_j=\frac{Cj+S+1}{\lambda j}.
\end{equation}
No independence of the $Z_v$ is assumed.
\end{proposition}
\begin{proof}
If the total excess exceeds $S$, its overloaded coordinates form a
nonempty subset $J$ whose total count is at least $C|J|+S+1$.
Markov's inequality with $z_j>1$, followed by a union over subsets of size
$j$, gives the expression. The zero-load case is deterministically safe.
\end{proof}

For main-only preparation the envelope follows from pointwise deletion.
At bit stage $b$, a key's raw destination depends only on its old source
block and the low $b$ bits of its new word. Earlier caps only remove keys,
so retained arrivals at every destination subset $J$ are at most the raw
arrivals of the complete logical population. Retention can depend on
earlier low bits of that same word; retained routes need not be independent.
Condition on the complete old-home counts $X_i$, and let $p_i(J)$ be the
probability that a complete key at old home $i$ has raw destination in $J$.
For $z\ge1$, the MGF of the complete raw count is
\[
 \E\!\left[z^{Z^{\rm raw}(J)}\mid(X_i)_i\right]
 =\prod_i(1+p_i(J)(z-1))^{X_i}.
\]
Uniform complete old homes, $N/m\le\lambda$, and destination column sums
$\sum_i p_i(J)\le |J|$ give the proposition's envelope. A proved joint
likelihood factor $D$ handles a nonuniform reference map. This is a positive
exponential comparison; it does not assert independent capped losses.

The finite sum is directly computable. If only subset sizes below $q$ are
evaluated, set $\mu=ne^\lambda/2^C$. Whenever $\mu/(q+1)<1$, the omitted
sizes are bounded by
\[
 D\,2^{-(S+1)}\frac{\mu^q}{q!}
              \frac{1}{1-\mu/(q+1)}.
\]
This follows by using $z=2$ and $\binom nj\le n^j/j!$. Otherwise increase
the cutoff or evaluate the remaining finite sum. The implementation is
\texttt{stage\_bound(n,S,capacity=C,mean=lambda)} in
\path{code/certify_preparation.py}, followed by the factor $D$.

\subsection{A sharper bound for completed populations}
\label{sec:positive-series}
\begin{proposition}[Joint retained-tier count]\label{prop:general-tiers}
Consider complete logical populations indexed by $i$, with at most $N_i$
distinct keys and $m_i$ main homes of capacity $C$. Assume their joint
home law is bounded by $D$ times a reference law in which the populations'
home-count vectors are independent uniform multinomials. Assume each
completed population has the full-home identity
\[
 T_i=\sum_{v=1}^{m_i}(X_{iv}-C)_+.
\]
Put $n=\sum_i m_i$ and $\lambda=\max_i N_i/m_i$. Any class of storage tests
whose physical inputs are subsets of these tiers, with no additional
sources and capacities at least $S$, has failure probability at most
$\min(1,D\Phi(n,C,\lambda,S))$, where $\Phi$ is computed below.
Shared keys in different tiers still count once per physical copy.
\end{proposition}
\begin{proof}
Any covered overflow implies $\sum_iT_i>S$. The indicator of this integer
event is bounded by the stop loss
$F_S(x)=(\sum_{i,v}(x_{iv}-C)_+-S)_+$.
This function is increasing and has nonnegative mixed differences for
nonnegative increments. For independent, identically distributed key
increments $V_1,\ldots,V_k$ and fixed $x$, telescoping gives
\[
 \E F_S(x+V_1+\cdots+V_k)
 \ge F_S(x)+k\{\E F_S(x+V_1)-F_S(x)\}.
\]
Averaging over $k\sim\Pois(1)$ and replacing keys one at a time increases
the expectation. Apply this to each independent complete population
separately. Poisson splitting gives independent comparison-home counts of
means $N_i/m_i$; pad those means to $\lambda$. This is a comparison of
expectations, not an assertion that the original homes are independent
Poisson variables. The joint density factor $D$ transfers the result.
\end{proof}

For $Z\sim\Pois(\lambda)$ define $w_j=\Prb(Z=C+j)$, $j\ge1$, and
$W(z)=\sum_{j\ge1}w_jz^j$. The generating function of $(Z-C)_+$ is
coefficientwise at most $1+W(z)$, so the comparison total has coefficients
bounded by
\[
 (1+W(z))^n\preceq\exp\{nW(z)\}=\sum_{k\ge0}c_kz^k,
 \qquad c_0=1,\quad c_k=\frac nk\sum_{j=1}^k j w_jc_{k-j}.
\]
The inequality uses $\binom nr\le n^r/r!$. With an integer cutoff $K>S$,
a computable upper bound on its stop loss is
\begin{equation}\label{eq:general-series}
 \Phi_K(n,C,\lambda,S)
 =\sum_{k=S+1}^K(k-S)c_k
 +\frac{K+1-S}{2^{K+1}}
       \exp\{n\,2^{-C}e^\lambda\}.
\end{equation}
Indeed $W(2)\le2^{-C}e^\lambda$ by dropping the event $Z>C$ in the
tilted Poisson expectation, and $(k-S)/2^k$ decreases for $k>K>S$.
The probabilities obey $p_0=e^{-\lambda}$ and
$p_k=\lambda p_{k-1}/k$, so all finite terms can be enclosed using positive
arithmetic. The implementation is
\texttt{positive\_excess\_bound(n,S,capacity=C,mean=lambda,cutoff=K)} in
\path{code/certify_tier_router.py}. Its current interface checks
$0<\lambda\le C$; the zero-load case needs no numerical routine.

Every integer cutoff $K>S$ gives a valid bound, so $K$ is an evaluation
parameter, not part of the statement. Unless a cutoff is passed, the
implementation takes the rule \texttt{series\_cutoff(n,S)}: the smallest
$K$ in $256,512,1024,\ldots$ with $K>S$ whose omitted-tail term, the last
term of Equation~\eqref{eq:general-series}, is at most $2^{-128}$. At
$C=64$, $\lambda=32$ and $S\le192$ this is 256 for every $n\le2^{24}$ homes, so every
value at those sizes is the cutoff-256 value; it is 512 at
$n=2^{25}$ and 1024 at $n=2^{26}$. At $n=2^{25}$ with 192 slots, where the
tail term dominated at cutoff 256, the top finish's merge/tier witness is
about $2^{-115.397}$ (\path{code/bound_table.json}) instead of the
$2^{-43.759}$ the same routine gives with \texttt{cutoff=256}.

Independent old/current families on jointly fixed complete key sets give
one application of the reference-law premise, even when those sets overlap.
Two selected tiers from the \emph{same} family do not automatically satisfy
it. Arbitrary valid source combinations therefore cannot be passed to this
formula without checking their dependence. In particular, selected route
sheds are not completed full-home tiers.

\section{Resolve before compression and top excess}
\begin{lemma}[Resolve before home compression]
\label{lem:resolved-compression}
Let an expanded reference representation and its residual lists contain
all versions to be merged. A public map groups its temporary homes into
final homes, and every copy of a key belongs to the same final home.
For each home, first resolve all versions by freshness, remove a winning
tombstone together with every version it shadows, then retain up to $C$
live winners and emit the rest. Assume no older source outside the merge
requires the tombstone. If $X_v$ counts the live winners at final home $v$,
then the final main count and total emitted demand are
\[
 M_v=\min(X_v,C),\qquad E=\sum_v(X_v-C)_+.
\]
Every append prefix of those emissions is at most $E$. A fused evaluator
with sufficient public local workspace and the same resolution-before-test
order has this same demand without materializing the expanded main table.
\end{lemma}
\begin{proof}
After resolution, home $v$ contains exactly $X_v$ distinct live keys.
Compaction retains $\min(X_v,C)$ and emits $(X_v-C)_+$. Summing gives
$E$, and an append prefix is a subset of these emissions. Fusing the
operations changes neither count when all relevant versions are available
before the capacity test.
\end{proof}
This is the intended short route for top finish. For a fixed logical trace,
the live winner set after consuming a carry is defined by logical history,
not by routing. Once its size and ideal-home premises are proved, the
existing subset/excess bounds apply directly to $E$. In the binary case,
a doubled representation can group two temporary homes into each final
home, with residual lists included in the resolution. A constant larger
fan-in changes the local workspace and map, not this counting argument.
The hierarchy fanout and the top's digit/way count are distinct parameters;
the native top sweep uses the latter.

This lemma does not bound the expanded representation's earlier routing
or its residual distribution. Nor does final-output equivalence alone prove
capacity equivalence: a fused evaluator must not overflow before performing
a deletion the reference performs first. A sweep that tests its shed before
offering B's tier, and drops tier tombstones only after the residual
epilogue, lets a tier tombstone consume capacity or resolve an older key only
after the sweep's test.

For example, at abstract capacity one, let A's main contain live key $a$,
B's main contain live key $b$, and B's tier contain a tombstone for $a$,
all targeting one final home. Both B records have the same carry home,
so its main/tier fullness identity holds. Such a sweep first sheds $b$;
the tier then replaces $a$ by its tombstone, which the final drop removes.
Only $b$ is a live winner, so the reference has $E=0$, while that sweep
emits one record. This is a local state satisfying the stated
placement/fullness conditions, not a claim about its frequency in engine
histories. The native top therefore resolves first.

\begin{proposition}[The native top finish resolves before it compresses]
\label{prop:top-resolution}
Let the top have $D$ ways, $1\le D\le8$, a carry of $U$ buckets, and final
homes $v=i+qU$ for carry bucket $i$ and way $q$. At the finish, C holds A's
main after preparation, R is A's retained carry, and each destination's
candidate pages are its leaf of R followed by B's tier T, which every
destination reads directly. Assume:
(P1) A (C and R) and B (main and T) are each distinct by key, and nothing
older than A and B is live;
(P2) every copy of a key has one final home $v$: a C record lies at $v$, a
B-main record at carry bucket $v\bmod U$ with way $\lfloor v/U\rfloor$, an
R record is a candidate naming $v$ on the leaf of $v$'s range, and a T
record is a candidate naming $v$ on B's tier pages;
(P3) the pool dedup has removed R's copy of every key T holds.
Let $W_v$ be the resolved live winners of home $v$ (freshest version of
each key, winning tombstones removed with what they shadow), $X_v=|W_v|$ and
$E=\sum_v(X_v-C)_+$. Then after the sweep $C[v]\subseteq W_v$ with
$|C[v]|=\min(X_v,C)$, and the sheds and leftovers naming $v$ are
$W_v\setminus C[v]$. Every append prefix of the merge stash and of the
output tier is at most $E$, the tier ends with exactly $E$ records, and
neither C nor the tier holds a tombstone.
\end{proposition}
\begin{proof}
Fix $v$, with carry bucket $i$ and way $q$. The visit of B's bucket $i$
first clears C's copies of all of B's main keys. For each of its
destinations, before the merge and over every page,
\texttt{ResolveTopDestination} (a) drops the R candidates for $v$ whose key
B's main holds in direction $q$, which by (P2) holds every B-main copy of a
key homed at $v$; (b) clears C[$v$]'s copies of the T candidates for $v$;
(c) drops every tombstone candidate; and (d) drops C[$v$]'s own tombstones.
The merge then omits B's main tombstones. By (P1)--(P3) each key homed at
$v$ now appears at most once among C[$v$], B's surviving direction-$q$
records and the candidates for $v$, and it appears exactly when its
freshest version is live: an A copy survives only if B holds no version
(B-main copies were cleared by the dedup and by (a), T copies by (b) and
(P3)); B's tombstones leave only after clearing A's copies; and an A
tombstone has nothing older (P1). So the survivors are $W_v$. The merge
fills free lanes and sheds only into a full C[$v$]; an offer fills free
lanes and leaves a candidate only when C[$v$] is full. After the resolution
nothing removes a record from C[$v$]: no candidate shares a key with a
record of C[$v$], a merged or shed B record or another candidate, so the
offers clear no lane, and no tombstone remains for the trailing drops.
Hence either C[$v$]$=W_v$ and nothing naming $v$ is emitted, or C[$v$] ends
full and exactly $W_v\setminus C[v]$ is emitted; summing gives $E$. The
merge stash receives sheds only. The epilogue drains the range leftovers,
runs the leftover dedup, which removes nothing (an older-class leftover
sharing a key with a shed would be an R candidate whose key B's main
holds), and then drains the merge stash; every prefix is a sublist of the
emitted records. The argument is per destination and $D$ enters only as the
direction a destination reads, so it covers every public $D$, including
$D=1$ and nonbinary maps.
\end{proof}
Lemma~\ref{lem:resolved-compression} therefore describes the native
finish's demand. The merge stash and the output tier are distinct buffers
bounded by the same witness $E$; R's DFS leaves are the finish
distributor's separate event (Corollary~\ref{cor:top-finish-separated}),
B's directly read tier pages have none, and the retained carry's capacity
is the preparation's. Proposition~\ref{prop:top-excess} supplies the
live-set size and home law for a symbolic probability bound on $E$.

\emph{Implementation.} \texttt{ResolveTopDestination}
(\path{src/oram/rebuild/cpu_byteslice_second_chance.hpp}) runs inside the
offer-fused fixed-merge sweep whenever that sweep may drop every tombstone
(\texttt{FixedMergeTombstones::kAll}), which is the level arena's top. It reads the
candidate pages in place, the A leaf as the older class and B's direct tier
pages as the fresher one, and clears only private masks, so it needs no scratch;
B stays read-only. A sweep unit is one source bucket with all of its
destinations, so pauses fall between complete resolutions, and lanes clear
disjoint candidates of AND-folded private range masks. Because no candidate
can then clear a lane, the offer allocates without its class-marked dedup;
the staler-first allocation that guards a staler copy of a shed key is
unchanged at every other site. The leftover dedup and the trailing drops
are kept as public work, and checked builds assert that the drops find
nothing. The resolution's candidate match and two intersects per page are
charged two units and the offer without dedup one: the pass weight plus one
per page.

\emph{Tests.} \path{tests/oram/cpu_byteslice_top_resolution_test.cpp}
drives the real sweep and epilogue on native-capacity inputs that satisfy
(P1)--(P3) and the fullness identities, against a scalar
resolve-then-compress oracle: the example above at capacity 64; a deleting
B-tier version on a later page than the candidates it makes room for;
$D=1,\ldots,8$ with one to three pages; B-main and B-tier tombstones, live
B-tier replacements and absent-key tombstones; mixed full and nonfull
homes; the merge stash and tier exactly full and one over; empty sources;
A tombstones, a premise violation that (d) handles; and slice-by-slice and
two-lane schedules. It checks every append prefix against $E$,
$\min(X_v,64)$ per home, the absence of tombstones before the drops, and
B's immutability. The job-level top battery checks every case against
the same identity under pauses, lanes and helpers, and adds a case in which
B's tier deletes a full bucket's keys while B's main adds live keys to it,
at $D=1,2,3,4,6,8$. A test-before-resolve sweep fails these fixtures (for
example, the capacity-64 example sheds one record with $E=0$). They are local
fixtures, not a reachability or frequency claim for engine histories.

\subsection{The top live population and its excess}
\label{sec:top-excess}
\begin{proposition}[One top-finish excess witness]
\label{prop:top-excess}
Let the public live capacity be $K$, the top have $T$ homes of capacity
$C$, and let $p_v$ be the probability of home $v$ under its current map.
Write $p_* = \max_v p_v$, $\lambda_*=Kp_*$ and
$\Gamma=(Tp_*)^K$. At each top ordinal, including ordinals not reached
following an earlier failure, define $W$ to be the live dictionary at its
logical window end and
\[
 X_v=\#\{x\in W:h_f(x)=v\},\qquad
 E_f=\sum_{v=0}^{T-1}(X_v-C)_+.
\]
If $S=\min(S_{\rm merge},S_{\rm tier})$, every first top-finish
merge-stash or tier-append overflow implies $E_f>S$. For
$0<\lambda_*\le C$ its probability is at most
\[
 b_{\rm top}(S)=\min\left\{1,
 \mathcal B_{\rm sub}(T,C,\lambda_*,S,1),
 \Gamma\Phi_{q}(T,C,K/T,S)\right\}.
\]
Here $q>S$ is any integer cutoff in Equation~\eqref{eq:general-series},
independent of the population capacity $K$.
For zero population the bound is zero. Charge this event once per top
ordinal, including all lane-local and fold tests at capacities at least
$S$; do not add another copy for the published top tier.
\end{proposition}
\begin{proof}
Proposition~\ref{prop:native-sources} gives a fixed set $W$ with $|W|\le K$.
Its home vector is multinomial with $|W|$ independent categorical trials
of probabilities $(p_v)$, even though W's physical realization and the
current-family B tier were computed with those same words. For any subset
$J$ and $z\ge1$,
\[
 \E z^{\sum_{v\in J}X_v}
 =\left(1+(z-1)\sum_{v\in J}p_v\right)^{|W|}
 \le\exp\{Kp_*|J|(z-1)\}.
\]
Proposition~\ref{prop:general-subset} gives the first probability bound.
Relative to uniform home assignments, the joint density of $|W|$ trials
is at most $(Tp_*)^{|W|}\le\Gamma$. Pad the uniform population to $K$
keys; excess is increasing. Proposition~\ref{prop:general-tiers} gives
the second bound for one complete population. Neither argument applies
a new independent law to selected survivors. Proposition~\ref{prop:top-resolution}
gives the pointwise append implication; the lane-storage lemma below
covers partitions of the same demand. Previous preparations, range tests
and publications may be assumed successful for that implication alone.
The probability above remains unconditional.
\end{proof}

For native $U=2^k$, $T=DU$ and $w=64$, the exact top maps are
\[
\begin{array}{ll}
 D\in\{1,2,4,8\}:&h(t)=t\bmod T,\quad p_*=1/T,\\[2pt]
 D\in\{3,5,6,7\}:&h(t)=(t\bmod U)+U((t\mathbin{\rm div}U)\bmod D),\\
 &p_*={\lceil2^{w-k}/D\rceil}/{2^w}.
\end{array}
\]
In particular $D=6$ uses a whole-digit modulus six, not the ordinary
level's odd-part decomposition. A home in way $a$ has probability
$\bigl(\lfloor2^{w-k}/D\rfloor+ind_{a<2^{w-k}\bmod D}\bigr)/2^w$.
These formulas require sufficient word width for the stated map. Native
$K=32T$ gives mean 32 for binary tops and $\lambda_*=32Tp_*$ otherwise.
One may replace $\Gamma$ by $1/(1-K\delta)$ when
$\delta=Tp_*-1$ and $K\delta<1$; the exact power does not require that
last inequality. The worked binary-top numerical evaluation below uses this bound at the
shared merge/output capacity; other top maps remain symbolic here.

\begin{corollary}[Top stage stashes for all supported maps]
\label{cor:top-stages}
Keep the top map, population and notation of
Proposition~\ref{prop:top-excess}. A preparation bit-stage stash of
capacity $S_b$ has the same excess-tail bound $b_{\rm top}(S_b)$; a
whole-digit radix-stage stash does also. Union these bounds over the
actual stage slots. This covers binary and whole-digit top preparation,
including $D=6$, independently of the separate range distributor.
\end{corollary}
\begin{proof}
Write the route as binary stages inside its public blocks, then the radix
stage if present. For a bit stage $b$ and old home $O_x$, the complete
raw destination is
\[
 V_b(x)=2^b\lfloor O_x/2^b\rfloor+(H_{f,x}\bmod2^b).
\]
Old and current families are independent on complete A. A destination's
raw probability is $2^{-b}$ times the old probability of its aligned
$2^b$-home block, hence at most $p_*$. Different keys' raw destinations
are independent. The retained main is a deletion of complete A; earlier
caps only delete further raw arrivals. Each stage stash therefore has
demand at most $\sum_v(X_{b,v}-C)_+$. The top-excess MGF and uniform-law
density arguments apply verbatim to this raw categorical law. At the
radix stage the complete raw destination is the final current home, with
probabilities $p_v$, giving the same conclusion. The final fused offers
run after the routing test of their disjoint pair or radix tuple; they
cannot increase another tuple's arrivals. Slots shared by the disjoint
blocks of one stage use one combined $T$-home witness, not one capacity
allowance per block.
\end{proof}

Initialization's scatter is fixed independently of the routing family.
If it puts at most $L_0\le C$ records per source home, a bit-stage raw
arrival has mean at most $L_0$ at each destination: each key chooses
uniformly within its aligned $2^b$ block. Distinct keys are independent,
though their categorical distributions differ. Their subset MGF is at
most $\exp\{L_0|J|(z-1)\}$. Marked Poisson replacement, exactly as in the
binary initialization proof, also gives the positive-series excess bound
$\Phi_q(T,C,L_0,S_b)$. Take the smaller of this and
$\mathcal B_{\rm sub}(T,C,L_0,S_b,1)$ for each bit-stage stash. If
$2L_0\le C$, the first bit stage is deterministically safe. The final
whole-digit radix stage, when present, has complete raw homes under the
current map and uses $b_{\rm top}(S_{\rm radix})$. Packing the initial
tier uses its input-excess bound from the same map. Thus initialization's
scatter, route stashes and publication-tier demand have symbolic bounds
for every native top map; only its distributor and the deterministic
native inventory remain separate obligations. No final-range theorem
for a nonbinary map is inferred, and the library refuses such tops.

\paragraph{Preparation, initialization and the distributor boundary.}
Preparation starts with the previous top's distinct live set A, of size
at most $K$, under an independent old family. Its retained carry has
exactly the new-home excess of A: A-only routing and offers conserve A,
and every residual targets a full home. Apply the same one-population
bound with the retained-carry capacity. Initialization starts with its
fixed distinct input set and publishes exactly its full-home excess in
the tier; charge that tier under initialization, not as a top-finish
exposure. The input scatter is deterministically safe (at most $C/2$ per
bucket); its intermediate routing stashes and distributor are separate.
For binary top preparation and initialization the stage and packing
bounds of Appendix~\ref{sec:binary-top-preparation} have their complete
source and early-selection premises from the native induction.

The finish's swept pool is R and B's carry tier, with combined
class-marked dedup before any distribution; the merge stash, which receives
the sweep's own sheds, is drain-only. Even a key
that B will delete at resolution can occupy R here. Thus $E_f$ does
\emph{not} bound the finish's distributor: the dictionary at the old cut can
be full while the new cut is empty. The distributor is bounded instead by
R's own law. R is the complete-home excess of the fixed set A, and a coarse
range sees that excess restricted to its homes; the following lemma is
that marginal.

\begin{lemma}[Complete-home excess on a public home set]
\label{lem:restricted-excess}
Let $W$ be a fixed set of at most $K$ keys whose homes $h(x)$ are
independent categorical trials on $T$ homes of capacity $C$, with
probabilities $p_v\le p_*$. For a public home set $J$ put
$Y_v=\#\{x\in W:h(x)=v\}$ and $R_J=\sum_{v\in J}(Y_v-C)_+$. For integers
$s\ge0$ and $q>s$, if $0<Kp_*\le C$,
\[
 \Prb\{R_J>s\}\le b_J(s)=\min\{1,\Phi_q(|J|,C,Kp_*,s)\}.
\]
An empty $J$ or $W$ gives zero.
\end{lemma}
\begin{proof}
Give a key whose home lies outside $J$ a zero increment; the keys'
increments to $(Y_v)_{v\in J}$ are then independent. The stop loss
$(R_J-s)_+$ is increasing and has nonnegative mixed differences for
nonnegative increments, so the marked Poisson comparison in the proof of
Proposition~\ref{prop:general-tiers} applies to this one complete
population: Poisson splitting gives independent counts on $J$ with means
$|W|p_v\le Kp_*$, and padding the means only increases the stop loss.
Equation~\eqref{eq:general-series} bounds the resulting expectation. No
uniformity of the $p_v$ within $J$ and no old-map density factor is used.
\end{proof}

\begin{corollary}[The top finish is a separated pass]
\label{cor:top-finish-separated}
Let the top be binary, with $T=2^b$ homes of capacity $C$ and $p_*=1/T$,
and let the public live capacity be $K\le CT/2$. At a top finish let R be
the preparation's retained carry, of physical capacity $S_R$, let B's
carry tier have $s_B$ pages, and let the merge stash and output tier have
capacities $S_{\rm merge}$ and $S_{\rm tier}$. Let the finish range plan
partition the homes into $2^\rho$ public coarse ranges
$J_r=\{v:\lfloor v/2^I\rfloor=r\}$ of $2^I$ homes, $\rho+I=b$. Suppose the
finish distributes R alone through one-page DFS leaves on this partition
with the capacities of Equation~\eqref{eq:dfs-capacity} ($S=C$,
$P_{\rm pool}=S_R$, $R_f=0$), and offers each destination its A leaf
followed by B's $s_B$ tier pages, read directly. Preserve the combined
class-marked dedup of R and B's tier, the
per-destination resolution of Proposition~\ref{prop:top-resolution}, and
the coordinated drain. Then every first capacity failure of the finish
implies one of
\begin{enumerate}
\item $E_A=\{\max_r R_r>C\}$, where $R_r=\sum_{v\in J_r}(Y^A_v-C)_+$ and
      $Y^A_v$ counts the previous top's live set A at current home $v$;
\item $E_f>\min(S_{\rm merge},S_{\rm tier})$, with $E_f$ as in
      Proposition~\ref{prop:top-excess};
\item a capacity failure of B's tier at its producer, charged there.
\end{enumerate}
A finish therefore contributes at most
\[
 b_{\rm fin}=\min\Bigl\{1,\;\ind[S_R>C]\,2^\rho\,\Phi_q(2^I,C,K/T,C)
   +b_{\rm top}\bigl(\min(S_{\rm merge},S_{\rm tier})\bigr)\Bigr\},
\]
and a direct plan, which routes nothing, has no first term. Neither the
top digit $D$ nor any A/B independence enters $b_{\rm fin}$.
\end{corollary}
\begin{proof}
This is Proposition~\ref{prop:separated-capacity} with A the retained
carry and B the carry tier; only its A marginal and its item~3 need their
top instances.
For item~1, A-only preparation leaves exactly $(Y^A_v-C)_+$ records of
home $v$ in R (the paragraph on preparation above). Combined dedup and the
unshadowed drops only delete them, so at most $R_r$ records of R are routed
to range $r$, and R holds at most $S_R$ records before the first failure.
Item~1 of the proposition (DFS containment) turns any first
A-distributor overflow into $E_A$, which is impossible when $S_R\le C$.
Lemma~\ref{lem:restricted-excess} with $W=A$, $J=J_r$ and $s=C$ gives
$\Prb\{R_r>C\}\le\Phi_q(2^I,C,K/T,C)$; take the union over the ranges.
Item~2 of the proposition holds as stated: every destination reads its one
A page and the same $s_B$ B pages, and B adds no distributor test.
For item~3, Proposition~\ref{prop:top-resolution} is per destination and
reads every candidate page of a destination, its A leaf and B's direct
pages alike, so it holds for this layout: every append prefix of the merge
stash and of the output tier is at most $E_f$, the count of
Lemma~\ref{lem:resolved-compression} for the resolved live set.
Proposition~\ref{prop:top-excess} bounds $\Prb\{E_f>S\}$ by
$b_{\rm top}(S)$, charged once for both buffers and every lane
(Lemma~\ref{lem:pause-storage}). Adding the two probabilities is
Theorem~\ref{thm:abstract-composition} for one finish. The top digit
enters only through the carry width $U=T/D$, which none of these
quantities reads.
\end{proof}

\paragraph{Worked values.}
At the worked shape $T=2^{22}$, R has two pages ($S_R=128$) and the finish
range plan has $\rho=4$ and $I=18$: 16 coarse ranges of $2^{18}$ homes.
The positive series (at the cutoff its rule picks,
Equation~\eqref{eq:general-series}, which is 256 at this shape) gives about
$2^{-82.775}$ per range, so about $2^{-78.775}$ per finish; the merge/output witness at capacity 128 is
about $2^{-126.056}$ per finish, charged once. Over the 257 public ordinals
of $D=8$, including speculation, the top finish contributes about
$2^{-70.769}$. With the same per-finish bound, $D=4,2,1$ have 129, 65 and
33 exposures and contribute about $2^{-71.764}$, $2^{-72.753}$ and
$2^{-73.731}$ (\path{code/configuration_space_2p27.json}).
At $T=2^{23}$, $2^{24}$ and $2^{25}$ ($N=2^{28}$ to $2^{30}$) the plan's
$2^{18}$-home width cap (the separated offer contract) gives $2^5$,
$2^6$ and $2^7$ ranges of $2^{18}$ homes, so the range term is about
$2^{-77.775}$, $2^{-76.775}$ and $2^{-75.775}$ per finish
(\path{code/bound_table.json}).
\path{code/evaluate_top_finish.py} evaluates the corollary at the native
row of each selected configuration: \path{code/configuration_probe.cpp}
prints the finish range plan (\texttt{LevelFinishRangePlan}), the direct B
pages and the pool slots from \texttt{TopMergePoolShapeOf}. The result is
\path{code/top_finish_evaluation.json}. The checker attaches it as three
obligations built by the same \texttt{separated\_finish\_obligations} as at
every sparse ordinary digit: the A-only range event, B's producer reuse
(charged zero) and the shared merge/tier witness. The selected ordinary
jobs plus this top finish contribute about $2^{-69.678}$ at both selected
ladders, but remaining source classes still prevent a full certificate.

 For \emph{top preparation and initialization}, the selected fix
is to reuse DFS at the existing proved binary-top coarse partition, with
one page per leaf; see Appendix~\ref{sec:top-dfs-decision}.
The finish uses the same distributor on another pool and partition; the
preparation's final-range bound is not reused there.

\begin{lemma}[Pauses, lane folds and retained storage]
\label{lem:pause-storage}
Suppose a logical stash accumulates a list of records with nonnegative
unit demand $D$ during a phase, with no removals before its fold, and has
capacity $S$. Split the list into any public lane partition. Give each
lane capacity at least $S$ and fold into an initially empty stash of
capacity $S$, checking every append. Then every lane overflow or fold
overflow implies $D>S$. The conclusion also holds with an existing
retained population if $D$ includes it and the fold appends to it exactly
once. Pausing and resuming without changing buffers or repeating operations
introduces no additional witness. An exact copy to capacity at least that
of its valid source is deterministically safe.
\end{lemma}
\begin{proof}
Every lane prefix and every fold prefix is a sublist of the combined
population counted by $D$. If $D\le S$, all those prefixes fit. Retention
and pauses leave the lists unchanged. A valid source has no more records
than its capacity, which proves the copy assertion.
\end{proof}
For a reused stash, $D$ includes all its preparation and cross writers;
bounding each writer separately by $S$ is insufficient. For DFS, the
public node capacities and full-parent work arrays satisfy the separate
containment proposition in Section~\ref{sec:dfs-theorem}. Retained buffers
need distinct reserved storage, or a proof that their lifetimes do not
overlap. This lemma does not bound total scratch allocation or allow two
paused jobs to overwrite each other's storage. Those are deterministic
ownership and allocation obligations. Lanes with smaller private capacities
would also need their own bounds.

\paragraph{Horizon accounting without a resumption factor.}
If potential admissions of a job class are at least $\tau$ access positions
apart, at most $\lceil H/\tau\rceil$ lie in $(0,H]$.
Indeed, $n$ such positions require $(n-1)\tau<H$.
If a public list of at most $J$ potential jobs admitted at or before the
left boundary covers every job that can still perform a tested operation
in the horizon, a whole-job witness has
at most
\[
 n(H)\le J+\lceil H/\tau\rceil
\]
exposures, plus any separately charged initialization. Enumerate these by
public logical ordinals, including jobs not reached after an earlier
failure. A partial job or a resumed job uses its original exposure. Cold
start has $J=0$ once initialization is separately handled; an arbitrary
warm start needs an unconditioned history bound for its $J$ jobs, not a
fresh-family assumption for their retained state.
For a fixed non-top digit the admission spacing is $b_\ell aP_\ell$ in
the logical counter schedule. A proof must still bound pre-boundary jobs
and show that speculative preparation belongs to the charged ordinal;
wall-clock starts and completions need not have this spacing. Bounding only
the number of randomly selected active jobs is insufficient: the public
potential-job list or a selection-safe witness is required.

\section{Public storage exposure accounting}
\label{sec:native-ledger}
Retained scratch and lane copies carry subsets of their producing witnesses.
They do not multiply event counts. Public extents must cover every copied
page, mask, task-list entry and DFS work array. Online dense storage is
bounded by admitted access positions rather than random occupancy.
\begin{proposition}[An ordinal exposure envelope including speculation]
\label{prop:ordinal-exposures}
Let $aP_\ell$ be the access positions per incoming level-$\ell$ carry.
Let $s_\ell$ be the seeded accepted-carry counter, zero for an empty
start, or the counter from \texttt{DirectInitializationSeed}. For $H$
online access positions after that start, set
\[
 m_\ell(H)=\left\lfloor\frac{H}{aP_\ell}\right\rfloor+1.
\]
Every non-initialization job that can perform record-storage work through
that horizon has ordinal in the public list
$\{s_\ell+1,\ldots,s_\ell+m_\ell(H)\}$.
For a non-top digit $d\in\{1,\ldots,b_\ell\}$ the number in that list is
\[
 n_{\ell,d}(H)=
 \left\lfloor\frac{s_\ell+m_\ell(H)-d}{b_\ell}\right\rfloor
 -\left\lfloor\frac{s_\ell-d}{b_\ell}\right\rfloor.
\]
Negative arguments here use mathematical floor. At the top there are
$m_L(H)$ potential ordinals, all at its fixed top digit. Initialization
adds one separate exposure if performed. Adoption-only digits have zero
job-demand bounds, even if included in this envelope.
\end{proposition}
\begin{proof}
A complete incoming carry requires its complete interval of $aP_\ell$
positions. At most $\lfloor H/(aP_\ell)\rfloor$ new such intervals exist.
FIFO and \texttt{Advance} bind only \texttt{ordinal()+1}; they cannot
prepare a second future ordinal before accepting the first. Speculation
can therefore add at most the next ordinal, regardless of thread speed,
pauses or helper interleaving. Counting ordinals congruent to $d$ modulo $b_\ell$ gives the formula.
The argument bounds potential jobs, rather than conditioning on which
jobs are active or complete at the right boundary.
\end{proof}
At an exact multiple of $aP_\ell$ the final $+1$ is intentional: a next
preparation can run before any of its carry's accesses. At a non-multiple
it equals the usual ceiling allowance. This envelope may overcount
unreachable speculative work, which is harmless. It assigns a whole-job
witness to partial jobs and never multiplies by the number of resumptions,
lanes, copies or retained scratch buffers.

For a horizon $(h_0,h_0+H]$ in the same initialized run, a completely
selection-safe warm-start envelope is the above list with $H$ replaced
by $h_0+H$, including its initialization if an unconditional bound on the
whole preceding run is desired. This conservatively charges all past
ordinals and covers every initially active or retained producer without
assuming fresh randomness for a snapshot. A shorter list consisting of
$J$ pre-boundary jobs plus new admissions is valid only after proving a
\emph{public} ordinal envelope for those $J$ jobs; a count of observed
active jobs is insufficient. Arbitrary imported warm states have no
claimed probability law here.

Consequently the symbolic record-storage sum can be written
\[
 B_{\rm records}(\Theta,H)=\min\left(1,
 b_{\rm init}+\sum_{\ell<L}\sum_{d=1}^{b_\ell}
 n_{\ell,d}(H)b_{\ell,d}
 +m_L(H)b_L\right),
\]
where each $b_{\ell,d}$ is the sum of its \emph{distinct} ledger witnesses,
and $b_L$ includes preparation, retained carry, both top distributors and
one finish-excess term. $b_{\rm init}=0$ for an empty start; otherwise it
includes scatter through publication. Open range terms receive bound one
until replaced by a proved applicable bound. Do not first maximize each
component over digits: sum within a digit, then apply its own exposure
count. This does not assert a complete engine certificate: each public shape requires fresh numerical applicability, a deterministic
inventory audit and discharge of the checker's premises.

\subsection{Structural applicability, separate from numerical certification}
\label{sec:structural-domain}
The window and exposure algebra permits arbitrary integer fanouts at least
two. The present native source-frontier proof and ordinary routing bounds
specialize to power-of-two $b_\ell\in\{2,4,8\}$ with binary carry widths,
$U_{\ell+1}=b_\ell U_\ell$ except for the explicitly modeled staged
low-load fold. Set $N_{A,\max}=(d-1)aP_\ell$,
$N_{B,\max}=aP_\ell$ and $\lambda_\ell=aP_\ell/U_\ell$ in ordinary
component bounds. Use the actual staged home width for a low-load route
and the actual carry width after its fold; do not silently substitute
$32U_\ell$ for both. Require the component's stated load, bit-width,
range-index and capacity inequalities. DFS permits public pages/widths
satisfying its containment layout; its source bound must use the matching
coarse ranges. Direct passes have no range-compression event.

The top-excess theorem covers the native maps for every $1\le D\le8$
allowed by geometry, with $T=DU_L$, any legal $K$, and the exact $p_*$
above. The binary-top stage/final-range adapter covers
binary tops only and retains its own size restrictions; the library
refuses other tops. The separated top finish
(Corollary~\ref{cor:top-finish-separated}) is stated for binary tops at
every $D$; its A marginal (Lemma~\ref{lem:restricted-excess}) needs only
$p_*$. Corollary~\ref{cor:top-stages} bounds every supported top
map's main-route stage stashes and initialization scatter; it does not
extend the final-range adapter to nonbinary tops. All theorems require
consistent family/key encoding without wraparound, successful public
allocation, sufficient copy/view extents, legal distinct initialization,
and a fixed logical trace obeying capacity at every prefix. Shared and
per-level tier profiles substitute their actual source, view and output
capacities; they do not alter the family argument.

A native geometry audit checks that a chosen implementation instance meets
these structural conditions and that every native test is assigned. A
numerical certificate evaluates its complete symbolic sum with rigorous
arithmetic and compares it with the requested target. Neither a successful
geometry probe nor a previously computed uniform-fanout subtotal performs
those other two tasks. Variable fanouts therefore need different public
sizes and exposure counts, not a new independence theorem. The default
half-fanout ladder is the first such ladder carried through all three
tasks; the configuration-space ledger lists, component by component, why
each premise holds at its fanout-four levels and at the four-to-eight
transition.

\section{Input invariants and probabilistic premises}
\subsection{Completed non-top tables}
At consuming digit $d>1$, A has $(d-1)U$ buckets and uses the \emph{previous}
digit's map and family. Its logical window has at most $32(d-1)U$ records,
including tombstones: a prefix contributes at most one record per access
position, and each subsequent union/dedup can only reduce the sum of its
input windows. Production geometry gives
$\texttt{PrefixPositions}\cdot\texttt{PrefixProduct}(\ell)=32U_\ell$.
Published main records are at their old homes, and the main/tier union is
distinct on a successful execution with valid inputs.

There is a stronger retained-load invariant than the arbitrary-split premise.
In a completed non-top table, every tier record targets a full main bucket.
An unabsorbed offer candidate has a full destination; a last fused-stage
shed also targets a bucket that just filled. Offers cannot lower occupancy:
every cleared main key has a distinct fresher replacement. Explicitly, with
main occupancy $m$, eligible stale/fresh counts $s,f$, and $q\le\min(m,f)$
cleared keys, the two rounds allocate
\[
 u=\min(s,64-m),\qquad v=\min(f,64-m+q-u).
\]
Then $m-q+u+v\ge m$, and it equals 64 whenever $s>u$ or $f>v$.
There is no subsequent non-top tombstone drop. Leftover dedup and drain
do not lower main occupancy. Adoption preserves the result, and
hashed main/tier publication makes both inputs non-removable by online
probes. Consequently,
\[
 a_i=\min(x_i,64).
\]
This is a count identity. The preparation proof below needs only
$a_i\le X_i$; neither symmetry of retained identities nor exchangeability
of selected loads is required.
The native job tests recompute every output tier key's home and check
that it is full, alongside their existing distinctness, placement and
newest-value checks, including sliced preparation and helper execution.
Checked and unchecked suites pass. They use smaller geometry; the separate
header probe enumerates the selected production geometry without tables.

\subsection{Exact odd-radix bias}
Uniform ideal words do not give exactly uniform homes for an odd-radix
old table. Write $M=r2^k$, $Q=\lfloor2^{64-k}/r\rfloor$ and
$s=2^{64-k}\bmod r$. A bucket in radix class $j$ has mass
\[
 p_j=\frac{Q+\ind_{j<s}}{2^{64-k}}\frac1{2^k}.
\]
Odd-radix boundaries have $r>1$. For a fixed logical set
of $N$ distinct keys with independent ideal old words, compare its old-home
law $P$ to uniform homes $P_0$. With $\delta=M\max_j p_j-1$,
\[
 \frac{dP}{dP_0}\le(1+\delta)^N\le\frac1{1-N\delta}\quad(N\delta<1).
\]
The last inequality expands the binomial and bounds its coefficients by
$N^j$. Applying the same fresh-mark/retention procedure to both laws
preserves this upper density for nonnegative tests. At $N\le32M$, the
largest audited factor is
\[
 D=\frac{67108864}{67108843}<1.000001.
\]
This multiplicative correction does not charge an additive modulo-bias
failure probability. Exact finite checks verify the pushed-forward bound
for retained count states and demonstrate their nonexchangeability under
biased homes. The comparison does not assign fresh laws to selected
residuals or condition ideal words on capacity success.

\subsection{Families and ownership}
The native family audit checks different families for A and the new merge,
and identical producer/consumer families for the carry. It also checks consecutive top families and the
reserved empty-top/initialization separation. These public identities
support the ideal-family model; Proposition~\ref{prop:native-sources}
supplies the complete-source induction on a fixed logical trace. Shared A/B keys use one common new
mark, and B's current-family selected tier is not a fresh independent source.

Published source payload stays pinned. Borrowed tier masks are copied into
private descriptors before modifying offers/deduplication; paused jobs
retain their cursor and scratch. Physical private copies do not add logical
keys. Tombstones still shadow older values and must retain their metadata.
Top preparation and initialization keep their existing separate paths and
explicit source/reset premises. Appendix~\ref{sec:native-ledger} charges partially completed jobs and
retained storage; unresolved range terms outside the selected
configurations and the checker's undischarged premises still prevent a
complete certificate (the deterministic native inventory requires a matching
configuration audit, Appendix~\ref{sec:deterministic-inventory}).

\section{Ordinary preparation: applying the subset bound}
\label{sec:preparation}
Before fixed B insertion, A's retained main count at
old home $i$ obeys $a_i\le X_i$. Its uncapped routes at bit stage $b$ lie
within disjoint binary source blocks: each source sends probability
$2^{-b}$ to each of its $2^b$ destinations. Thus every destination column
has source-probability sum at most one. With $N/M\le32$, deletion gives
exactly the envelope of Proposition~\ref{prop:general-subset}, with
$n\le(d-1)U_\ell$, $C=64$, $\lambda=32$ and old-map factor $D$.
Earlier capping only removes arrivals; no sheds return before finishing.
The bound needs neither independent stages nor exchangeable retained loads.

Use the allocated stash capacity $S\in\{64,128\}$ and sum over at most
\texttt{block\_bits} stages. The numerical evaluator retains subset sizes
through 64 and uses the general tail formula with cutoff $q=65$.

This covers A's preparation population only. A's old tier is excluded,
and later cross writers sharing a stash require the combined-demand proof
in the per-key Poissonization comparison.

\subsection{Final fused stashes and both output-tier drain phases}
\label{sec:final-drain-bound}
There is a simpler joint witness for the storage \emph{after} the range
router. Let $K=A\cup B$ be the complete distinct logical union, including
tombstones, with $|K|\le32dU$. Assume this complete set is fixed independently
of its new ideal-family words; shared keys have one word, and distinct keys
have independent words. This assumption concerns the complete source, not
its selected tier or a law conditioned on earlier capacity success. Define
\[
 X_v=\#\{x\in K:D(H_x)=v\},\qquad
 E=\sum_{v=0}^{dU-1}(X_v-64)_+.
\]
No uniform old-home law is required for this witness.

\paragraph{The final insertion stash.}
In fixed insertion, meeting A/B copies deduplicate before the capacity
test and retagging selects the exact binary or radix final way.
A itself has no internal duplicates.
Consequently each final destination receives a distinct subset of the raw
$X_v$ keys. Its shed count is at most $(X_v-64)_+$. The insertion drain-only
stash has no preparation writer and is separate from every A stage stash.
It starts empty and only accumulates these final sheds. Every accumulation
prefix is therefore at most $E$, even if the next append would fail.
Offers do not add records to this stash. This accounts for one shared stash
capacity, without giving each source a separate allowance.

\paragraph{The output tier before and after leftover dedup.}
For the argument, extend only the tier's capacity to infinity and follow
the same execution until its first would-be overflow. The swept pool
deduplicates by source rank before scatter; its range leftovers therefore
contain at most one physical copy of each key. After all offers, each
leftover is disjoint from final main and targets a full home. At a full
home there are only $X_v-64$ distinct union keys absent from main, so the
first drain of range leftovers, including each packing prefix, is bounded
by $E$. This is true \emph{before} deduplicating against the drain-only
stash; no temporary second tier capacity is assumed.

The drain-only epilogue must also inspect the \emph{original} fresh swept
sources, including candidates already absorbed into main. For example, a
fresh candidate can replace a main key while a different fresh candidate
occupies the freed slot and duplicates a final-stage shed. Reading only
the remaining range candidates would miss that shed/main duplicate. The
native pool dedup reads the original stash and clears the shed copy.
Any remaining cross-set duplicates are resolved by leftover dedup after
the first drain and before drain-only records append. Both sides are then
distinct from each other and from main, and every residual has a full home.
All subsequent packing prefixes are again at most $E$. With conservation
of the complete union and no non-top tombstone drop, the final main/tier
counts are exactly $\min(X_v,64)$ and $(X_v-64)_+$ at each home.

This proof covers direct level-one digit-one output drainage. Higher-level
digit-one adoption reuses the producer tier witness at its actual capacity.

\paragraph{Current-map probability.}
Under uniform final homes, for any destination subset $J$ and $z\ge1$,
\[
 \E z^{\sum_{v\in J}X_v}
 =\left(1+\frac{|J|}{dU}(z-1)\right)^{|K|}
 \le \exp(32|J|(z-1)).
\]
Apply the subset witness of Equation~\eqref{eq:general-subset} with
$n=dU$ and the actual tested capacity $S$. For the current map's modulo
bias, use $\delta=dU\,p_{\max}-1$ and multiply once by
$1/(1-32dU\delta)$, as in the earlier old-map density argument. This
factor applies to the complete union's new homes; it does not assert
freshness or independence for selected A/B residuals.

Finite tests check 6,000 shared-key
offer/drain histories and small exact multinomial tails.

\section{Binary-top preparation and initialization}
The top rekeys A before B insertion. Its old tier and all preparation
losses form an A-only pool; initialization has distinct input keys and
an empty old tier. One-page DFS distributes that pool. The final-range
witness below contains every DFS node overflow.
\subsection{A shadow process counts each shed key once}
To avoid assuming independent losses at successive stages, retain an
analytical shadow of each key that leaves main. At a bucket, the cap sends
\[
 (M,R)\longmapsto(\min(M,64),R+(M-64)_+),
\]
where $M$ is main occupancy and $R$ counts already-shed shadows. Every key
is either main or shadow. Both follow their unused routing bits, and no
shed returns to main during the counted interval. Survivor selection must
ignore the bits not yet used. This is a premise about the input flow,
not a consequence of distinct keys alone.

The terminal test $(\sum R-S)_+$ is increasing with increasing differences,
and a shadow increment has at least the effect of a main increment.
Backward expectation through routing preserves these properties: an added
record chooses one output by an independent unused bit, and main/shadow
increments can be coupled to that same choice. The cap preserves them as
well. Below capacity an added main record stays main; at capacity its
effect changes to a shadow increment, which is larger. Thus the boundary
second difference remains nonnegative, as do the other mixed differences.

Suppose the first counted arrivals are subsets of independent-key ideal
positions with mean at most $\lam$ per coordinate. Monotonicity permits
replacement by the complete ideal positions. The per-key Poissonization
argument of the per-key Poissonization comparison, applied to the backward test,
then gives independent $\Pois(\lam)$ initial populations after padding.
This is an expectation comparison after averaging unused bits; it does
not assert pathwise domination by independent capped buckets. Establishing
the ideal-position subset premise for each application remains necessary.

\subsection{The finite loss tree and its error bound}
At depth one draw $N\sim\Pois(\lam)$ and set
$(M,R)=(\min(N,64),(N-64)_+)$. For the next depth take two independent
child states. Thin each child's main and shadow populations independently
with probability $1/2$, giving combined counts $U,V$, and set
\[
 (M',R')=(\min(U,64),V+(U-64)_+).
\]
After depth $t$, thin $R$ by $2^{-h}$ for $h$ still-unrevealed pattern bits;
call the result $L_{t,h}$. This recurrence keeps all losses along one
counted route together, rather than multiplying successive-stage MGFs.

The distribution is enclosed by exact nonnegative integer arithmetic in
units of $2^{-512}$, with cutoff $H=256$. At every tree node the ideal
main-plus-shadow population is $\Pois(\lam)$. Write
\[
 \tau_H\le\frac{\Prb\{\Pois(\lam)=H+1\}}{1-\lam/(H+2)},
 \qquad \mu_h=\lam 2^{-h}.
\]
A depth-$t$ tree has $2^t-1$ nodes. The thinned loss $L_{t,h}$ is at most
the thinned root main-plus-shadow total, whose law is $\Pois(\mu_h)$;
therefore $\E L_{t,h}^2\le\mu_h+\mu_h^2$. A union bound on a truncated
node and Cauchy--Schwarz give, for $g$ roots,
the additive stop-loss correction
\[
 \epsilon=g\sqrt{(\mu_h+\mu_h^2)(2^t-1)\tau_H}.
\]
For each root, replace its loss by zero on that root's truncation event.
The positive-series convention $c_0=1$ majorizes this completed law:
it retains all positive untruncated masses and puts at least the missing
mass at zero. Thus convolution does not discard other roots' contributions
when one root truncates. The positive-part test is 1-Lipschitz, so its
error is at most the sum of $\E[L_{t,h};\text{root truncates}]$;
adding an independent early population does not increase this correction.
This justifies both the single-root and sorted-top calculations.
Rounding and truncation are included in the numerical bounds below.

\subsection{Parameterized binary-top stage and packing bounds}
\label{sec:binary-top-preparation}
The binary-top adapter covers $T=2^k$ homes, $k\ge2$, main capacity 64,
one stage-stash capacity $S$, and an output drain capacity $R$.
For preparation, $R=S$ is the retained preparation carry; for
initialization, $R$ is the published tier capacity. The implementation still
has a separate old tier of capacity $R_{\rm old}$. Require one distinct,
tombstone-free complete A population of at most $32T$ keys, independent
of the relevant ideal families. These premises concern A before the
incoming B merge; no half-load assertion for the later A/B union is inferred.
The stage and packing series use any cutoff $q>\max(S,R)$, selected
adaptively by \texttt{series\_cutoff}. The loss-tree cutoff $K=256$ is
separate: its final-range test is at 64 slots, while $S$ and $R_{\rm old}$
enter only the early binomial population cap $IS+R_{\rm old}$ (or $IS$
for initialization). That binomial law is not capped at $K$: its finite
coefficients and its full generating function at 2 bound the omitted
series. The tree's truncation correction remains valid after adding the
early population. Thus $S,R\ge256$ pose no proof restriction and require
no larger loss-tree cutoff. At $N=2^{31}$ a stage-stash width of 3 buckets
does not make the bound small: $\Phi$ per stage is about
$2^{-71.017}$ before unioning stages and exposures. The existing series
suffices with stage-stash width 4, giving about $2^{-119.717}$ per stage,
and the top's single packing event fits the lifetime budget with its
three-page tier. The sizing header takes this as a floor for any level
whose output exceeds $2^{25}$ homes, keeping a larger horizon-based
requirement if present; middle-level stage stashes at $D=8$ keep their
existing widths. This is the smallest whole-page width that closes the
term with the positive-series bound. This capacity choice is justified by
the evaluated proof, not by the loss-tree cutoff; the routing and the tree
are unaffected.

\paragraph{Raw stage arrivals.}
In preparation, let $O_x$ be the old binary home and $H_x$ the independent
new ideal word of key $x$. At one-based stage $b$, its uncapped destination is
\[
 V_b(x)=2^b\left\lfloor O_x/2^b\right\rfloor+(H_x\bmod2^b).
\]
Uniform independent old homes and new words make $V_b(x)$ uniform on
$[0,T)$, independently for distinct keys. Main membership may depend on
the old family: the retained main is still a subset of this complete raw
population, and earlier route deletions only remove keys. Earlier stages
perform no offers; the final fused stage offers only after routing each
disjoint pair. Its offers cannot change another pair's routing input.
The stage stash demand is therefore pointwise at most
\[
 E_b=\sum_{v=0}^{T-1}(X_{b,v}-64)_+,
 \qquad X_{b,v}=\#\{x:V_b(x)=v\}.
\]
This uses deletion, not exchangeability of selected retained loads.

Initialization instead has a public, family-independent source scatter
with at most 32 keys per bucket. For stage $b$, each key chooses uniformly
among $2^b$ destinations in its fixed source block. These independent
categorical laws need not be identical across keys. Every destination's
mean is at most 32. Apply the increasing-differences marked-Poisson
comparison from the per-key Poissonization comparison to
$(\sum_v(X_{b,v}-64)_+-S)_+$: replace each independent key by an independent
$\Pois(1)$ number of copies with independent marks. Poisson splitting
then gives independent destination counts at their respective means;
pad each mean to 32. The same comparison applies to preparation's
multinomial raw counts. Independence is introduced in this comparison,
not assumed of capped bucket populations or different stages.

Write $Q(T,s)$ for the positive-series upper bound computed by
\texttt{positive\_excess\_bound(T,s)} at capacity 64 and mean 32.
The conditional stage-stash union bounds are
\[
 kQ(T,S)\quad\hbox{for preparation},\qquad
 (k-1)Q(T,S)\quad\hbox{for initialization}.
\]
The missing first initialization term is deterministic: two initial buckets
contain at most 64 records together, so neither first-stage destination
can exceed 64. No independence between stage failures is needed.

\paragraph{One A-only packing witness.}
Let $X_v$ count the complete A population at its new final home, and
$E=\sum_v(X_v-64)_+$. The old tier, earlier sheds, retained main and last
stage's shed partition one distinct population. The actual top preparation
pass uses \texttt{kNoDedup}. A successful range distributor preserves that
partition; final offers only insert into free lanes. Every range leftover
therefore targets a full home and is disjoint from main. Every final-stage
shed also targets a home that became full before its offer and stays full.
Thus their combined population is at most $E$ (equal to $E$ if all keys
are conserved). Draining the range leftovers first and the disjoint
drain-only stash second gives the same upper bound on every packing prefix.
There is no temporary second drain capacity and no early deduplication of
physical copies. This establishes $Q(T,R)$ for the retained preparation
carry or initial published tier. Range-distribution failures remain separate
events; the reference-word law is not conditioned on their absence.

\paragraph{Parameterized sorted ranges.}
For range index width $I$, assume $0\le I<k-1$, one scatter group with
eight slots per initial range bucket, and $2^{k-I}$ final ranges. The final
main stage $k$ is drain-only. Early stages $1,\ldots,I$ and the old tier
give a deferred binomial padding population
\[
 P=IS+R_{\rm old}\quad\hbox{for preparation},\qquad
 P_0=IS\quad\hbox{for initialization},\qquad p=2^{-(k-I)}.
\]
Later stages $I+1,\ldots,k-1$ give depth $t=k-I-1$, one remaining thinning
bit and $g=2^{I+1}$ shadow roots. Substituting these parameters into the
following sorted-range derivation generalizes its positive-series and
truncation calculation. Disjoint-key deferred selection and the ideal-arrival
reset comparison remain explicit premises. The native source and frontier inductions above supply these histories
for binary preparation and fixed-scatter initialization; the existing
shadow argument supplies the reset comparison.

\paragraph{Computation and checks.}
\path{code/certify_top_preparation.py} applies these formulas to eight
distinct physical/init shapes obtained from the 16 native domain settings
(each default-ladder setting ends in the same top as its uniform partner);
\path{code/top_preparation_certificates.json} records the arithmetic,
source hashes and tree caches. At $T=2^{22}$ and $S=R=128$, preparation's
stage union is below $2^{-121}$ and its packing bound is below $2^{-126}$
per pass. Adding sorted final ranges gives a conditional contribution below
$2^{-108}$ over 257 regular preparation allowances plus one initialization
(displayed logarithm $-108.1670$). This does not include top finish.

Seven finite tests check raw categorical laws, a correlated-family
counterexample, initial scatter, 2,000 A-only offer/drain histories, the
generalized sorted formula against the previous fixed-size certificate,
and applicability. Forty-one native fixtures check final-stage and tier
64/65 and 128/129 boundaries, earlier-pool offers and actual initialization
scatter. These adversarial kernel fixtures validate pointwise counts;
they are neither ideal-family probability samples nor a reachability proof.
Nonbinary tops (refused by the library) and the later B merge remain
outside this numerical adapter. The symbolic top results above do not
extend its final-range arithmetic.

\subsection{Final ranges of binary top preparation and initialization}
\label{sec:top-final-ranges}
This bound is stated for a sorted reference distributor (an analysis model,
not library code); it concerns only the swept pool and its final
partition, and Corollary~\ref{cor:top-dfs-transfer} transfers it to the
one-page DFS the library runs.
There are $2^{22}$ main buckets, 512 final ranges, and $I=13$ low bits
matched by the eventual offer. The range pattern uses bits 13 through 21.
Stages 1--21 enter the swept pool; stage 22 is drain-only. Each allocated
stage stash has 128 slots.

Separate stages 1--13 and the old tier from the later sheds. Their keys
have been selected without revealing their nine range bits, and are
disjoint from the remaining main keys. For a fixed range, expose the
remaining main-key words as well. The late loss is now determined while
the early count is still $\Bin(m,2^{-9})$, where
\[
 m\le P=14\cdot128\quad\text{for preparation},\qquad
 m\le P_0=13\cdot128\quad\text{for initialization}.
\]
Padding to $P$ or $P_0$ yields an independent binomial comparison for
increasing tests. This does not assume that the actual early and late
counts were independent. It relies on disjoint keys and unused-bit-blind
selection, supplied by the A-only frontier in
Proposition~\ref{prop:native-frontiers}; shared copies would invalidate it.

Stages 14--21 give $t=8$, $h=1$, and $g=2^{14}$ comparison roots for one
range. Consequently its overflow probability is bounded by
\[
 \E\left(B+\sum_{i=1}^{g}L_{8,1}^{(i)}-64\right)_+,
 \qquad B\sim\Bin(P,2^{-9}),
\]
with independent quantities only in the comparison experiment. Let $w_j$
enclose the positive masses of the truncated root loss and set
\[
 C(z)=(1-p+pz)^P\exp\left(g\sum_{j=1}^{H}w_jz^j\right),\qquad p=2^{-9}.
\]
Use the positive coefficient recurrence of Section~\ref{sec:positive-series}
and convolve with the exact binomial coefficients. The series tail is at
most $(H+1-64)2^{-(H+1)}(1+p)^P\exp(gW(2))$; add $\epsilon$ for root
truncation. Union over the 512 final ranges.

This gives $<2^{-117}$ per initialization, $<2^{-116}$ per top preparation,
and $<2^{-108}$ after charging the larger bound for all 257 exposures,
including initialization. Earlier stage failures are counted separately;
the ideal experiment is not conditioned on earlier success. These numbers
do not bound the top's later B merge (the finish).

\subsection{Selected implementation: one-page DFS for top preparation}
\label{sec:top-dfs-decision}
The preparation builder always uses the existing DFS distributor (the
geometry refuses a nonbinary top). Initialization uses the same builder. It preserves the swept source pool and coarse partition of the
binary-top final-range proof, with one 64-slot page per leaf.
The top finish is a separate pass with its own pool and partition: it
distributes only the retained carry, whose law is A's complete-home excess
(Corollary~\ref{cor:top-finish-separated}), and this final-range bound is
not used there.
One page minimizes the page count and offer multiplicity at this fixed
partition within the present interface. No measured global runtime optimum,
universal one-page bound, or wider-range setting is claimed.

\paragraph{Freshness and the actual distinction between levels.}
A's old tier was selected under the old family. Relative to that selection,
its destinations under the new family are fresh, at the top just as at
ordinary levels. The newly produced stage sheds were selected using some
of the \emph{current} family's bits; only their unexamined bits remain
fresh. The existing shadow argument handles these dependent losses when
bounding a final range. This is not a special loss of freshness at the top.
The implementation difference is that ordinary preparation defers its
residual offers to the ordinary finish, where DFS is wired,
whereas the top fully routes and offers A in its own preparation pass,
which always distributes through its own one-page DFS. Initialization uses
that same preparation path. This integration requires no additional
source-independence premise.

\begin{corollary}[Transfer of the top final-range bound to DFS]
\label{cor:top-dfs-transfer}
Fix an A-only preparation or initialization instance satisfying the
binary-top source, stage-capacity and coarse-partition premises above.
Let $Y_r$ count its swept pool records naming final range $r$, and let
$\varepsilon_{\rm range}$ be the existing unconditional first-failure
bound for $\max_r Y_r>64$. Replace only its distributor by DFS, preserving
that pool and partition, with node capacities
\[
 C(h)=\min(P,64\cdot2^h),
\]
where $P$ is the public pool slot capacity and a height-$h$ node covers
$2^h$ leaves. Then the probability that this distributor is the first
capacity failure is at most $\varepsilon_{\rm range}$. Charge one range
event per preparation or initialization, with no factor for DFS nodes.
Main-route stash and retained-carry/publication-tier events remain separate.
\end{corollary}
\begin{proof}
On an execution with no earlier failure, the swept multiset is unchanged
by the replacement. The sorted reference's final-range test counts exactly these
$Y_r$; their population bound does not depend on sorting the records.
If every $Y_r\le64$, each ancestor's count is at most 64 times its number
of descendant leaves, and also at most $P$. The gather and work arrays
hold their full inputs, so the DFS containment theorem covers every
partition capacity test. Hence a first DFS failure implies the same
reference range event used by the existing bound. The first-failure
implication does not condition the probability law on earlier success.
\end{proof}

\paragraph{One page here does not improve the ordinary bound.}
Appendix~\ref{sec:top-final-ranges} already bounded 64-slot final ranges
for the worked binary top: $T=2^{22}$, $I=13$, 512 coarse ranges and
128-slot stage stashes/old tier. DFS with one page asks for the \emph{same} event, so this
transfer uses no stronger probability estimate. A published top tier is yet another capacity:
one temporary page per range does not justify one bucket in that tier.

\paragraph{Native applicability and acceptance scope.}
\texttt{TopPrepareRangePlan} is the common per-pass plan for allocation,
work prediction and execution. It fixes one page and the DFS range plan's
coarse width, which is the production plan's (its $2^{18}$-home cap,
the separated offer contract, does not bind at any top preparation
from $2^{20}$ to $2^{30}$); the library has no page or width setting, and its geometry
refuses nonbinary tops. \texttt{MaxLevelDfsAllocation} and the widest range-table calculation
include this pass; engine scratch construction and byte accounting already
consume those helpers. The preparation and finish use the same exclusive
per-level scratch at disjoint times. The existing coarse-task reservation
covers gather and partition;
DFS units remain sliceable through the existing cursor.

The source pool, last-stage drain-only stash, final-stage fusion, private masks,
original-source lifetime and both drain phases are unchanged. Native top tests
check logical contents against the test model for binary digits
$1,2,4,8$, including retained excess, poisoned input tags, empty A, pauses,
lanes and helpers. Resumed execution's charged residual work equals the
planner estimate. Direct initialization tests cover empty, half and full
input for all four digits and read back contents through online rebuilds.
The existing distributor suite checks record bytes, masks, placement, empty
and full pools and the 64/65 boundary.

The checker matches all source capacities and partition dimensions to the
parameterized final-range certificate, then requires an exact native DFS row
in that certificate's audited shape list. Its scatter field is explicitly 64
for DFS instead of the production builder's 8: this is a build-format change,
not a population or coarse-partition change. It charges the range witness once,
separately from main-route and drain excess. Nonbinary maps, changed widths,
missing evidence and stale dependencies cannot inherit this bound.

\paragraph{Workspace.}
At a fixed range width, one DFS page retains one offer page per main
bucket. For the worked 44-source-bucket, 512-range shape, the
workspace formulas give 9,600 additional record slots plus 6,216 bytes
of masks, prefixes and unit metadata: about 278 KiB at 29-byte stored
records, excluding allocator rounding. The 512-bucket leaf table serves as
the range table. The DFS schedule has 512 units covering 1,008
compactions. Wider ranges may reduce work but require reevaluating their
population bound.

\section{Completed-A comparison}
\label{sec:separated-residuals}\label{sec:dfs-application}
\subsection{Completed-A comparison by padding the blank input}
\label{sec:completed-a-shadow}
The selected route admits a single shadow comparison, including its final
losses. It is unnecessary to keep separate preparation-span marginals or
allocate leaf capacity among successive A stages. The following statement
is symbolic; it does not reuse separate preparation-only numbers.

\begin{proposition}[Completed-A coarse-range bound]
Let the bucket capacity be a positive integer $C$, $\lambda>0$, $U=2^k$, $d=r2^e>1$ with $r$ odd,
$b=k+e<w$, and $T=dU=r2^b$. Assume the native ordinary A source contract:
a fixed distinct logical population of at most $\lambda(d-1)U$ keys,
old homes uniform on $(d-1)U$ positions up to a joint density factor $D$,
and an independent current $w$-bit ideal word per key. Retained main is a
subset of those old-home positions. Embed it at $[U,T)$, leave $[0,U)$
empty, and run the implemented binary blocks and terminal radix without
reinsertion or selection on unused current bits. Include the old tier and
\emph{all} A routing losses in the distributor population.

Choose $0\le I<k$. Let $S_j$ bound the entire shared stash for binary
stage $j$, across all its preparation and continuation writers, and let
$T_A$ bound the old tier. Put
\[
 P=T_A+\sum_{j=1}^{I}S_j,\quad g=2^I,\quad t=b-I,
 \qquad q_a=\frac{\#\{0\le h<2^{w-b}:h\bmod r=a\}}{2^{w-b}}.
\]
The early range probability is $p_a=q_a2^{-(b-I)}$.
Define $F_t=(M_t,R_t)$ by the existing binary shadow recurrence at capacity
$C$ and mean $\lambda$, with \emph{no final thinning}. For $r=1$ let
$Z_a=R_t$. For $r>1$, take $r$ independent copies of $F_t$ and,
conditionally on their states, independent counts
\[
 U_j\sim\operatorname{Bin}(M_t^{(j)},q_a),\qquad
 V_j\sim\operatorname{Bin}(R_t^{(j)},q_a),\qquad
 Z_a=\sum_{j=1}^r V_j+\left(\sum_{j=1}^r U_j-C\right)_+.
\]
All independence in these definitions belongs to the comparison process.
Let $E_a\sim\operatorname{Bin}(P,p_a)$ be independent of $g$ copies of
$Z_a$, and set $Y_a=E_a+\sum_{v=1}^{g}Z_a^{(v)}$. For nonnegative integral leaf
capacity $L$, the first-failure contribution of any coarse range in radix
class $a$ is at most
\[
 \min\{1,D\,\mathbb E[(Y_a-L)_+]\}.
\]
Consequently a whole completed-A distributor is bounded by
\[
 a_{\rm job}\le
 \min\left\{1,D\,2^{b-I}\sum_{a=0}^{r-1}
                   \mathbb E[(Y_a-L)_+]\right\}.                 \tag{CA}
\]
Every DFS storage overflow is contained in this event. In the selected
implementation $C=L=64$. This statement excludes top
preparation/finish, and widths $I\ge k$; a direct small-pool plan has no
leaf-capacity event and needs no use of (CA).
\end{proposition}

\begin{proof}
Use an analytical frontier after stages $1,\ldots,I$ in every binary
subnetwork. The physical schedule may finish disjoint blocks in a different
order; their operations commute. Early sheds and the old tier are distinct
and disjoint from the remaining A main. Their unused current bits retain
the deferred-mark law. On a history with no earlier capacity failure their
number is at most $P$. Apply the existing early-padding argument first:
replace their range count by an independent $\operatorname{Bin}(P,p_a)$.
Use the bounded reference list on other histories as in the earlier
first-failure argument. In particular, do not condition on B's history,
on the eventual pool deduplication, or on execution success.

Only then restore complete-A positions as phantom main records at the
first counted arrivals, stage $I+1$. In the uniform old-home reference law,
each occupied old position has mean at most $\lambda$. Both endpoints of
the blank interval are aligned to $2^{I+1}$, since $I<k$. Thus every first
arrival coordinate has mean at most $\lambda$: its old-address block is
wholly blank or wholly occupied, and its low current bits are uniform.
The per-key ideal-position comparison for the backward shadow test gives
independent Poisson coordinates of these means. Add independent Poisson
records to raise every coordinate, including the blank region, to
$\lambda$. This only increases the increasing shadow test. Padding is an
analytical replacement, not a change to the native route or its load.

The actual preparation spans followed by the new continuation levels are
exactly the full binary network in each of the $r$ blocks, with the blank
region's ineffective early operations omitted. For a fixed final range,
there are $g$ low-bit roots per block, each with $t$ counted cap stages.
The existing backward increasing-differences argument and Poisson splitting
give the binary shadow comparison $F_t$ jointly at these roots. This step
retains dependence between losses at successive stages.

If $r>1$, the untouched high word independently selects radix $a$ with
probability $q_a$ for each key. Previously shed shadows follow that same
selection; arriving main records undergo the terminal capacity-$C$ cap.
The displayed definition of $Z_a$ is precisely that operation. Its cap is
the same $(M,R)\mapsto(\min(M,C),R+(M-C)_+)$ map used at binary nodes;
independent thinning and addition preserve the same backward test class.
Therefore the earlier shadow comparison also applies with this terminal
node. No assumption of independent binary and radix losses is introduced.
For $r=1$ there is no terminal node, and $R_t$ already includes the last
binary cap. Apply the old-home density factor $D$ once to the full
nonnegative test. Later combined dedup can only delete A candidates.
Finally, an integer count above $L$ contributes at least one to its
stop-loss, and each radix class has $2^{b-I}$ coarse ranges. Union over
these ranges and apply DFS containment.
\end{proof}

\paragraph{Evaluation and remaining scope.}
For a joint-state probability generating function
$F(x,y)=\mathbb E[x^M y^R]$, the terminal node is particularly simple:
form
\[
 G(x,y)=F(1-q_a+q_a x,\,1-q_a+q_a y)^r,
\]
then move its coefficient at $(u,v)$ to
$(\min(u,C),v+(u-C)_+)$. The marginal of the second coordinate gives
$Z_a$. Keeping the joint $(M,R)$ law is essential: thinning an existing
loss-only marginal omits the newly shed main records. Replacing $q_a$ by
$1/r$ is also invalid for a finite-word biased radix; its exact value, or a
proved conservative replacement, is required.

\paragraph{A rigorous terminal truncation enclosure.}
Let $H>C$ be the count cutoff, and let
\[
 \tau_H(\nu)=\frac{\Prb\{\Pois(\nu)=H+1\}}{1-\nu/(H+2)},
 \qquad \nu<H+2.
\]
In each binary tree discard a state whenever its main-plus-shadow total
exceeds $H$, as in the existing binary calculation. At the terminal radix
node thin both coordinates with the exact $q_a$, convolve the $r$ children,
and discard totals above $H$ before applying the final cap. Intermediate
partial sums in that convolution may also be discarded above $H$:
nonnegative remaining summands cannot reduce their total.

Put $\mu_a=\lambda$ for $r=1$ and $\mu_a=r\lambda q_a$ otherwise. Every
binary node has total $\Pois(\lambda)$, while the terminal total is
$\Pois(\mu_a)$: caps preserve totals, and independent thinning and addition
preserve the Poisson law. A root's bad-event probability is at most
\[
 \beta_a=r(2^t-1)\tau_H(\lambda)
              +\ind_{\{r>1\}}\tau_H(\mu_a).
\]
The discarded partial sums require no extra union term because their
failure implies the terminal-total event. Since the root loss $Z_a$ is
at most that root's total, Cauchy--Schwarz gives
\[
 \E[Z_a;\text{some node discarded}]
       \le\sqrt{(\mu_a+\mu_a^2)\beta_a}.
\]
Replace a root's loss by zero on this event. The stop-loss function is
1-Lipschitz, so for $g$ roots, including the independent early population,
the additive correction per coarse range is
\[
 \boxed{\epsilon_a=g\sqrt{(\mu_a+\mu_a^2)
       \left[r(2^t-1)\tau_H(\lambda)
                 +\ind_{\{r>1\}}\tau_H(\mu_a)\right]}.}
\]
This specializes to the existing unthinned binary correction when $r=1$.
It requires no independence between truncation and the terminal total.

\paragraph{Positive arithmetic and the final stop-loss.}
The driver \path{code/certify_completed_a.py} uses $H=256$ and
$2^{-512}$ fixed-point upper coefficients. It computes each binomial
thinning weight from the exact rational $q_a$, rounds upward, retains the
joint law through all radix additions, and caps only after the last
addition. Packed integer convolution uses a coefficient-width bound
calculated from its actual inputs, so there are no carries between
coefficients. The optional GMP accelerator changes only execution time.
No floating-point convolution is used.

Let $w_j$ majorize the retained root loss mass at $j>0$, and put
$W(z)=\sum_{j=1}^H w_jz^j$ and $B(z)=(1-p_a+p_az)^P$.
The zero-on-failure completed root law is coefficientwise at most
$1+W(z)$; hence $B(z)\exp(gW(z))$ majorizes the completed sum law.
Writing $c_n=[z^n]B(z)\exp(gW(z))$, the driver bounds its stop-loss by
\[
 \sum_{n=L+1}^H(n-L)c_n+
 \frac{H+1-L}{2^{H+1}}B(2)\exp(gW(2)).
\]
Here $(n-L)2^{-n}$ is nonincreasing for $n>H>L$. The coefficients through
$H$ use the positive exponential recurrence; the displayed remainder,
Poisson tails, square root and final sums are all rounded upward.
Add $\epsilon_a$ before multiplying by the density and range count in
(CA). This also retains other roots' contributions when one root is
truncated; simply convolving deficient subprobability laws would not.

\begin{corollary}[Complete raw excess for separated middle levels]
\label{cor:middle-storage}
Let $M=(d-1)U$, $T=dU=r2^b$, with $U=2^k$, odd $r$ and
$2\le d\le8$. Assume the native source law of
Proposition~\ref{prop:native-sources} and the separated ordinary route
and drain premises. Write $D_A$ for the old-A home-law density bound,
$D_{A,\rm current}$ for A's current-home density bound and
$D_{\cup,\rm current}$ for the complete distinct union's current-home
density bound. Put
\[
 G(T,S)=\min\{\mathcal B_{\rm sub}(T,64,32,S,1),
                         \Phi_q(T,64,32,S)\},\qquad q>S.
\]
The A binary stage stashes of capacity $S_A$ have total failure bound
$\min(1,bD_A G(T,S_A))$. The terminal radix stash, when present, has
bound $\min(1,D_{A,\rm current}G(T,S_A))$. One event bounds both the
insertion stash and every published-tier append prefix by
\[
 \min\{1,D_{\cup,\rm current}
       G(T,\min(S_{\rm insert},S_{\rm output}))\}.
\]
These are unconditional reference-event bounds for first-failure
composition, for every middle level with the stated shape.
\end{corollary}
\begin{proof}
The needed additional premise beyond a subset-MGF envelope is that the
\emph{complete raw destinations}, under the reference law, are independent
categorical trials across distinct keys. Proposition~\ref{prop:native-sources}
already supplies it: complete logical sets are fixed independently of
routing words, and old and current families have their joint independent
law. This premise is not asserted for selected main survivors.

At binary stage $j$, a complete A key with old home $O\in[0,M)$ and
current word $H$ has uncapped destination
\[
 V_j=2^j\left\lfloor\frac{O+U}{2^j}\right\rfloor+(H\bmod2^j).
\]
Under independent uniform old homes, an aligned block contains at most
$2^j$ admitted old positions. Consequently every destination has
probability at most $1/M$, and the destinations of distinct keys are
independent. They need not be uniform on the expanded $T$ homes.
There are at most $32M$ keys. Apply Lemma~\ref{lem:restricted-excess}
to this categorical population with $J=[0,T)$, $p_*=1/M$ and population
cap $32M$: its complete raw excess has tail at most
$\Phi_q(T,64,32,S_A)$. The same home-mass bound gives the subset-MGF
envelope at mean 32, hence Proposition~\ref{prop:general-subset} gives
$\mathcal B_{\rm sub}$ as well. Transfer both bounds by $D_A$ once.
This is the marked-Poisson comparison already proved in
Proposition~\ref{prop:general-tiers}, in the nonuniform form of
Lemma~\ref{lem:restricted-excess}; no new comparison theorem is needed.

Retained A records and earlier caps only delete complete raw arrivals.
Preparation and continuation writers of one binary stage visit disjoint
nodes of the expanded blank-B network, once each. Their combined shared
stash demand is therefore at most the single complete raw excess.
Union over its $b$ stages; empty and non-shedding operations may be
included. At the terminal radix, the complete raw destinations are A's
current final homes. Use the uniform current-home reference, at most
$32M\le32T$ keys, and transfer by $D_{A,\rm current}$ once.

For insertion and publication, the existing complete-home witness
(Proposition~\ref{prop:separated-capacity} and the coordinated-drain
argument) uses the fixed distinct logical A/B union, of size at most
$32T$. Shared keys have one current word and are counted once. Its
reference homes are independent uniform choices on $T$; apply
Proposition~\ref{prop:general-tiers} to this one population, or
Lemma~\ref{lem:restricted-excess} with $p_*=1/T$, and the subset bound.
Transfer by $D_{\cup,\rm current}$, not by the old-map factor: insertion
and publication depend on current union homes. Charge once at the smaller
capacity. Neither the two inputs' physical dependence nor conditioning
on earlier success enters the probability calculation.
\end{proof}

After the completed A route, the fixed B insertion deduplicates meeting
keys before its capacity test, chooses each key's exact final way, and
then offers residuals. Its total sheds are at most the complete distinct
A/B union's current-home excess $E$. The existing coordinated-drain proof
bounds every output-tier append prefix by the same $E$, including the
range-only drain before final shed deduplication. Thus one event
\[
 E>\min(S_{\rm insert},S_{\rm output})
\]
covers insertion and publication. The native pool recipe gives insertion
its own initially empty drain-only stash, separate from all A binary and
radix stashes. B's already published tier is charged at its producer, with
no new consumer, lane or range multiplier. These are ordinary contracts:
tombstones remain logical records; top live-set resolution is separate.

\paragraph{Native applicability.}
The separated-source adapter checks every A bit/terminal stash once, the old
tier, direct B-tier pages, the initially empty insertion stash and actual
one-page DFS plan. The completed-A probability comparison applies at direct
level-one shapes as well: its input is the completed A population, not B's
source representation. Require $0\le I<k$ and exact public-row matching.
Missing or stale evidence is charged one. Each digit's stage, range and
shared insertion/output witnesses are summed before multiplying by its
public ordinal count.

\section{Published-tier sizing and the per-level profile}
\label{sec:tier-sizing}
For a completed non-top population with final-home counts $X_v$, the
published tier has the full-home identity
\[
 T=\sum_v(X_v-64)_+.
\]
For any population retaining the same final home map, this sum is at least
the overflow storage required after filling main, regardless of whether
its distributor uses sorting or a butterfly. Better distribution does not
make an overfull home larger.

\paragraph{A direct finite-population lower witness at the top.}
The admitted top capacity is $N=32m$ distinct live records with $m=2^{22}$
homes; initialization explicitly permits this population. Under the ideal
uniform-home law the occupancy vector is multinomial. For a one-bucket
tier, the event $\sum_v(X_v-64)_+=65$ already forces failure. We bound that
event \emph{from below in the finite multinomial law}, rather than treating
a Poisson upper comparison as a lower bound.

More generally let $C$ be main capacity, $S$ tier capacity, $t=S+1$,
$\lambda=N/m\le C/2$, and $m>t+1$. With
\[
 W(z)=\sum_{e\ge1}\Prb\{\Pois(\lambda)=C+e\}z^e,
 \qquad c_t=[z^t]\exp(mW(z)),
\]
consider $j\le t$ specified overloaded homes with positive excesses summing
to $t$. Their total record count is $K=Cj+t$. Summing their multinomial
probabilities over home choices and excess compositions differs from the
corresponding contribution to $c_t$ by the factor
\[
 \frac{(m)_j}{m^j}\frac{(N)_K}{N^K}
           e^{\lambda j}(1-j/m)^{N-K},
\]
where $(n)_k$ denotes a falling factorial. Let $K_*=(C+1)t$. The product
inequality $\prod_i(1-x_i)\ge1-\sum_i x_i$, and
$\log(1-x)\ge-x-x^2/(2(1-x))$, give the uniform lower factor
\[
 a=\left(1-\frac{t(t-1)}{2m}\right)
   \left(1-\frac{K_*(K_*-1)}{2N}\right)
   \left(1-\frac{\lambda t^2}{2(m-t)}\right),
\]
provided all three factors are positive. For the last factor, discard the
positive $Kj/m$ term and use $e^{-x}\ge1-x$.

Conditioned on those specified exact occupancies, the remaining
$n=N-K$ keys choose among $m'=m-j$ homes uniformly. For every count $u$,
\[
 \Prb\{\Bin(n,1/m')=u\}
 \le e^{n/(m'(m'-1))}
         \Prb\{\Pois(n/(m'-1))=u\}.
\]
This follows directly from $\binom nu\le n^u/u!$ and
$(1-1/m')^n\le e^{-n/m'}$. Since $C\ge2\lambda$, we have
$K\ge\lambda(j+1)$ and hence $n/(m'-1)\le\lambda$. A conditional union
bound therefore bounds the probability of \emph{any additional} overloaded
home by
\[
 b=m\exp\!\left(\frac{\lambda}{m-t-1}\right)
                  \Prb\{\Pois(\lambda)>C\}.
\]
When $b<1$, demand that all other homes have at most $C$ records. The
events for different sets of overloaded homes and compositions are then
disjoint, yielding the finite-population lower bound
\[
 \Prb\!\left\{\sum_v(X_v-C)_+=t\right\}\ge a(1-b)c_t.
\]
The code evaluates $c_t$ downward using the positive recurrence, evaluates
$b$ upward with an explicit Poisson-series tail, and performs the final
comparison as an exact rational inequality.

For $m=2^{22}$, $N=32m$, $C=S=64$, the result is
\[
 \Prb\{\text{required tier size}>64\}>2^{-62}>2^{-64},
 \qquad \text{displayed lower-bound }\log_2\approx-61.8572.
\]
This is a single-population obstruction, already above the entire lifetime
target. It applies to the declared full live initialization population,
not to every smaller workload or an estimated measured failure rate.
Earlier algorithmic failures do not evade the storage requirement.
Consequently bitonic sorting cannot justify reducing \emph{every}
published tier to one bucket at the current maximum population and home
map. A smaller population/load or a different storage design is a separate
configuration question.

The analysis and exact small-multinomial checks are in
\path{code/analyze_tier_sizes.py}, \path{code/test_tier_sizes.py}, and
\path{code/tier_size_analysis.json}.

\paragraph{The per-level profile.} The engine sizes each level's tier by
its own lifetime requirement for the tables it produces, a Chernoff bound
over that level's homes and executions, kept non-decreasing below the top;
this is the automatic sizing. A tier's width is its producer's and never
changes over time: level $\ell$'s banks and top banks hold its own width
and level $\ell+1$'s inboxes hold it too, so a merge reads its carry at the
level below's width and its older operand at its own, except an adopted
first carry, still in its inbox at its producer's width; forwards keep one
width on both sides. Each job's bound therefore depends only on its input
tiers' widths and its own, which are fields of its native shape key, so
the per-job table of Appendix~\ref{sec:configuration-space} covers the
profile with no new statement. At the default horizon the profile is
1,1,1,1,1,2 at $2^{22}$ entries and 1,1,1,1,1,2,2 on the default ladder at
$2^{27}$. It leaves the total bytes almost unchanged but narrows every lower
level's pool and the tier each of its tables adds to a probe, so a probe
touches fewer tier buckets than with one shared width. The
engine checks each profile against the sizing, and the engine tests run
pinned profiles of every seam's width change against a reference model.
A pinned shared width is a different shape and is not covered by the
artifacts.


\section{Deterministic inventory and current ledger}
\label{sec:deterministic-inventory}\label{prop:deterministic-inventory}
\path{code/audit_deterministic_inventory.py} enumerates tested capacities and
fixed extents in the engine include closure. A matching native plan assigns
every reachable test to a component witness or a deterministic capacity.
For direct level one the assignment is:
\begin{center}\small
\begin{tabular}{p{0.37\linewidth}p{0.54\linewidth}}
\toprule
Native test & Reference event or deterministic reason\\\midrule
Digit-one dense drain & B complete-home excess at output-tier capacity\\
Dense drain at $d>1$ & B complete-home excess at one producer bucket\\
Completed-A shared stage stash & Complete raw A stage excess\\
A-only DFS node/leaf & One completed-A coarse-range event\\
Fixed insertion stash and lane copies & Dominated by output capacity;
complete distinct-union excess\\
Output append/pack prefixes & Same complete distinct-union excess\\
Source/home-cache and mask views & Public indexed extents and copied masks\\
\bottomrule
\end{tabular}
\end{center}
Higher ordinary levels reuse completed-A separation. Binary-top preparation
and initialization use their stage, carry/tier and final-range witnesses.
Top finish uses retained-carry range excess and one resolved-live-set
insertion/output event. No imported B tier receives a consumer multiplier.

The audit's reviewed rules and lexical lint do not establish all semantic
claims that a planner reads public state only. The inventory is evidence
under those recorded assumptions. Liveness, allocations and memory safety
outside the enumerated extents remain outside its claim. An absent or stale
configuration section leaves its inventory line charged one.

\section{Configuration-space composition}
\label{sec:configuration-space}\label{prop:shape-key}
A bound-table key is its role plus its complete public native row, omitting
only the level index used in labels. It includes actual producer, older and
output tier widths, stash capacity, digit, routing partition and direct/DFS
mode. Direct level one additionally records the producer mode and standalone
one-bucket B tier. Every key must match the adapter's exact source contract.
Exposures depend on the configuration and are not part of the per-job bound.

The planner space fixes $B=8$, $q=2$, $H=2^{32}$ and automatic per-level tier
sizing, enumerates supported depths, $r_0\in\{2,4,8\}$, binary top digits
and native half-fanout ladders, and considers $s=1,2,4$ at capacities
$2^{15}$ through $2^{31}$. Larger capacities, other horizons, bank-first
placement, pinned tiers and explicit non-planner ladders need matching
applicability and inventory artifacts. A formula's validity alone does not
establish a target comparison.

The configuration-space sweep uses 401 of the bound table's 413 public
shapes. Across the 17 capacities,
334 of 339 candidates meet the conditional target at $s=1$, 332 of 337 at
$s=2$, and 328 of 333 at $s=4$. Every planner default meets the target,
and all spaces have complete native/inventory coverage with no missing job
keys. All candidates through $2^{30}$ meet the target; at $2^{31}$ five
$D=1$ geometries in each setting exceed it. This comparison remains
conditional on each candidate's recorded premises.

The actual A-only coarse width still must satisfy the completed-A theorem's
$0\le I<k$ condition. Native direct A-pool plans have no distributor event.
A larger prefix changes $U=2r_0s$ and the producer's $\Phi$ argument, and
changes the A route's public shape; it is evaluated rather than inferred
from the default prefix result.

\path{code/bound_table.json}, the selected configuration ledgers and
\path{code/configuration_space_2p}$n$\path{.json} (with the corresponding
$s=2,4$ filenames) contain exact outward-rounded rational bounds.
\path{CONFIGURATION_SPACE.md} presents their numerical outcomes.
\path{code/export_certified_bounds.py} exports conditional target results
with exact geometry to the planner's table. The library label ``certified''
means a matching conditional capacity bound is at most $2^{-64}$; it does
not discharge the checker's mathematical premises. Stale dependencies,
missing keys and refusals fail closed. Displayed logarithms are summaries;
acceptance uses exact rational comparisons.

\section{Open applicability and reproduction}
\label{sec:evidence}\label{sec:remaining}
The checker retains status \texttt{unproved\_applicability} and unconditional
upper bound one while native model/tape correspondence, trace and unshadowed
promises, imported-comparison applicability and inventory-rule semantics
remain open. A native geometry test or finite numerical check cannot replace
these arguments. The complete direct-level-one reference and the ideal
hierarchy probability chain are proved, with the scope stated in
Section~\ref{sec:lean-current}. \path{LEAN_CLAIMS.md} states the trusted
boundary and premises; \path{lean/README.md} explains how to build and audit
the proofs.

\path{code/README.md} maps tools to certificates and describes
regeneration in dependency order. Source/input hashes are regenerated by
their producers, rather than edited to accept stale computations. The direct
producer, paused jobs, lanes, helpers, freshness and output-placement native
regressions validate the implementation's concrete contracts. The completed-A
and abstract-completion finite checks validate recurrences and numerical
bounds without discharging the ideal-source laws.
\path{COMPLETED_A_FRONTIER.md} proves the abstract deferred-mark lemma
under explicit measurability premises. Its finite regressions distinguish
intermediate unused suffixes from completed residual selection; the latter
does not preserve an independent-family law.

From the repository root, run:
\begin{verbatim}
python3 analysis/bound_certification/code/regenerate.py --dry-run
python3 analysis/bound_certification/code/regenerate.py
\end{verbatim}
For report-only edits, from this directory run:
\begin{verbatim}
python3 code/build_report.py
python3 code/manifest.py
python3 code/manifest.py --check
\end{verbatim}
This report, the artifacts in \path{code/} and the registered Lean statements
define the analysis.
\end{document}
